\documentclass[10pt,a4paper]{article}

% ----------------- Paquetes -------------------%
\usepackage[T1]{fontenc}
\usepackage[utf8]{inputenc}
\usepackage[a4paper,margin=2.5cm]{geometry}
\usepackage{mathptmx}             
\usepackage{changepage} 
\usepackage{setspace}
\usepackage[spanish]{babel}
\usepackage[labelfont=bf,labelsep=space]{caption}
\usepackage{amssymb}
\usepackage{amsmath}
\usepackage{ebgaramond}
\usepackage{placeins}
\usepackage{etoolbox}
\usepackage{setspace}
\usepackage{parskip} 
\usepackage{tcolorbox}
\usepackage{needspace}
\usepackage{xparse}
\usepackage{ocgx2}
\usepackage{hyperref}
\usepackage{hypcap}
\usepackage{soul}\sethlcolor{yellow}% resaltado amarillo: \hl{texto} (Panel TCEE, ⌃H)


%%%%%%%%%%%%%%%%%%%%%%%%%%%%%%%%%%%%%%%%%%%%%%%%%%%%%%%%%%%%%%%%
% --- Fija primero tus valores globales "buenos" ---
\setstretch{1.2}                 % Interlineado global
\setlength{\parskip}{0.7\baselineskip}
\setlength{\parindent}{0pt}
\newlength{\GlobalParskip}
\setlength{\GlobalParskip}{\parskip}

\newcounter{eqnum}[subsection]
\renewcommand{\theeqnum}{\arabic{eqnum}}
\newcounter{alphaitem}
\newcounter{numitem}

%% ------------------- Recuadros----------------------%%
\newcommand{\puntosclave}[1]{%
  \begin{tcolorbox}[
    colframe=black,
    title={Puntos clave},
    before upper={\setlength{\parindent}{0.5cm}\setlength{\parskip}{0.5\baselineskip}}
  ]
  #1
  \end{tcolorbox}
}

\newcommand{\modificaciones}[1]{
  \begin{tcolorbox}[
    colback=white,
    colframe=red,
    title={Modificaciones a realizar},
    before upper={\setlength{\parindent}{0.5cm}\setlength{\parskip}{0.5\baselineskip}}
  ]
  #1
  \end{tcolorbox}
}

%% ------------------- Preguntas TEST----------------------%%
% Opciones: radio (single-choice)
\NewDocumentCommand{\OptRadio}{m m m}{%
  \noindent\CheckBox[radio,name=#1,value=#2,width=1.2ex,height=1.2ex]{}%
  \;\textbf{#2.}~#3\par
}

% Opciones: checkbox (multi-choice)
\NewDocumentCommand{\OptCheck}{m m}{%
  \noindent\CheckBox[name=#1,width=1.2ex,height=1.2ex,bordercolor={0 0 0}]{}%
  \;#2\par
}

% Pregunta single-choice (radio)
% Args: 1=id, 2=enunciado, 3=opciones, 4=solución+justificación, 5=imagen (o {}), 6=origen
\NewDocumentCommand{\PreguntaTestSingle}{m m m m m m}{%
  \par\medskip
  \noindent\textbf{#1.}~#2\par\smallskip

    % Imagen (si no es {})
    \IfBlankTF{#5}{}{%
      \begin{center}
        \includegraphics[width=0.75\linewidth]{#5}
      \end{center}
    }

    \begin{Form}
      #3
    \end{Form}

    \par\smallskip
    \switchocg{sol-#1}{\fbox{Mostrar / ocultar solución}}%
    \hfill{\footnotesize #6}\par\smallskip

    \begin{ocg}{Solución}{sol-#1}{0}
      \noindent #4
    \end{ocg}
  \par\medskip
}

% Pregunta multi-choice (checkboxes)
\NewDocumentCommand{\PreguntaTestMulti}{m m m m m m}{%
  \par\medskip
  \noindent\textbf{#1.}~#2\par\smallskip

    % Imagen (si no es {})
    \IfBlankTF{#5}{}{%
      \begin{center}
        \includegraphics[width=0.75\linewidth]{#5}
      \end{center}
    }
    \begin{Form}
      #3
    \end{Form}

    \par\smallskip
    \switchocg{sol-#1}{\fbox{Mostrar / ocultar solución}}%
    \hfill{\footnotesize #6}\par\smallskip

    \begin{ocg}{Solución}{sol-#1}{0}
      \noindent #4
    \end{ocg}
  \par\medskip
}


%% ------------------- CITA ----------------------%%
% \begin{cita}[Autor][Año][Obra] texto \end{cita}: cita sangrada y, debajo a la derecha, «AUTOR (Año), Obra». Los tres datos son opcionales.
% En el Panel TCEE: escribir \cita e Intro (el cursor va al texto; Tab pasa a Autor, Año y Obra).
\NewDocumentEnvironment{cita}{ O{} O{} O{} +b }{%
  \begin{adjustwidth}{2cm}{0pt}
  \raggedright
  #4\par
  \raggedleft
  \if\relax\detokenize{#1#2#3}\relax
    % sin firma
  \else
    #1%
    \if\relax\detokenize{#2}\relax\else\if\relax\detokenize{#1}\relax\else\space\fi(#2)\fi
    \if\relax\detokenize{#3}\relax\else\if\relax\detokenize{#1#2}\relax\else, \fi\textit{#3}\fi
  \fi
  \end{adjustwidth}
}{}



%% ------------------- Listas----------------------%%
    % --- Lista alfabética ---
    \newenvironment{la}
      {\begin{list}{\alph{alphaitem})}{%
          \usecounter{alphaitem}%
          \setlength{\leftmargin}{1cm}%
          \setlength{\parskip}{3pt}% 
          \setlength{\itemsep}{0pt}% ← espacio entre ítems
          \setlength{\parsep}{0pt}%  ← párrafos dentro de un ítem
      }}
      {\end{list}}
      
    % --- Lista numérica ---
    \newenvironment{lnum}
      {\begin{list}{\arabic{numitem}.}{%
          \usecounter{numitem}%
          \setlength{\leftmargin}{1cm}%
          \setlength{\parskip}{3pt}% 
          \setlength{\itemsep}{0pt}% ← espacio entre ítems
          \setlength{\parsep}{0pt}%  ← párrafos dentro de un ítem
      }}
      {\end{list}}
      
% ------------------- Imágenes----------------------%
    \graphicspath{{Img/}}% Rutas por defecto 
    % --- Imagen adaptativa ---
    % Registros de dimensión definidos una sola vez
    \newdimen\imgcap
    \newdimen\imgmin
    \newdimen\imgmax
    \newdimen\imgremain
    \newdimen\imgh
    \newdimen\imgneed
    
    \newcommand{\imagenfit}[3]{%
      \addvspace{10pt}% separación antes
      \par\begingroup
        % Parámetros de composición (asignaciones, no \newdimen)
        \imgcap=1\baselineskip            % alto estimado de título+fuente
        \imgmin=0.15\textheight
        \imgmax=0.15\textheight
        \imgremain=\pagegoal
        \ifdim\pagegoal=\maxdimen \imgremain=0pt\fi
        \advance\imgremain by -\pagetotal
        \advance\imgremain by -\imgcap
        % Decisión de altura destino
        \ifdim\imgremain<\imgmin
          \newpage
          \imgh=0.20\textheight
        \else
          \imgh=\imgremain
          \ifdim\imgh>\imgmax \imgh=\imgmax \fi
        \fi
        % Reservar espacio para evitar cortes del bloque completo
        \imgneed=\imgh \advance\imgneed by \imgcap
        \Needspace{\imgneed}
        % Cuerpo
        \noindent\begin{minipage}{\textwidth}
          \centering
          \captionof{figure}{\textemdash\ #2}\par\vspace{0.3em}% título arriba
          \includegraphics[width=\linewidth,height=\imgh,keepaspectratio]{\detokenize{#1}}\par
          \vspace{0.3em}\small\textit{Fuente: }#3% fuente abajo
        \end{minipage}%
      \endgroup
      \par\addvspace{10pt}%
    }

%------------ Bloque de RECUADRO ESPECIAL -------------%
    \newenvironment{specialblock}[1]{%
      \begin{quote} % crea sangría a izquierda y derecha
      \small % texto más pequeño
      \noindent\textbf{#1}\\[-0.3em] % título en negrita
      \rule{\linewidth}{0.4pt}\\[0.6em]}{
      \end{quote}}
      
%-------------------------- Portada -----------------------%
\newcommand{\oldtitle}[1]{\def\OldTitle{#1}}
\newcommand{\changerationale}[1]{\def\ChangeWhy{#1}}

\newcommand{\changebox}{%
\begin{tcolorbox}[title={Historial de cambios del tema}]
\textbf{Título anterior:} \OldTitle\par
\textbf{Motivo del cambio:} \ChangeWhy
\end{tcolorbox}
}

%-------------------------- Author Font -----------------------%
    \newcommand{\authorfont}[1]{{\fontfamily{EBGaramond-TLF}\selectfont #1}}

%-------------------------- Anexos -----------------------%
    \makeatletter
    \newcounter{anexo}
    \renewcommand{\theanexo}{\Roman{anexo}}
    \newcommand{\anexo}[1]{%
      \clearpage
      \phantomsection
      \refstepcounter{anexo}%
      \noindent\textbf{Anexo \theanexo.- }#1\par\medskip
      \addcontentsline{stc}{section}{Anexo \theanexo.- #1}%
    }
    \makeatother

    %-------------------------- Lista de figuras (personal) -----------------------%
\makeatletter
\newcommand{\MyListOfFigures}{%
  \clearpage
  \phantomsection
  {\centering\bfseries\Large Índice de figuras\par}%
  \medskip
  % Entrada en nuestro índice de contenido (stc)
  \addcontentsline{stc}{section}{Índice de figuras}%
  % Recuperación del listado (sin cabecera propia)
  \begingroup
    \parindent=0pt\parskip=2pt
    \@starttoc{lof}%
  \endgroup
}
\makeatother

%-------------------------- Bibliografía -----------------------%
\makeatletter
% Contador para las referencias
\newcounter{myref}
% Cabecera de la bibliografía, entra en el índice 'stc'
\newcommand{\MyBibliography}{%
  \clearpage
  \phantomsection
  {\centering\bfseries\Large Bibliografía y referencias\par}%
  \medskip
  \addcontentsline{stc}{section}{Bibliografía y referencias}%
  \setcounter{myref}{0}%
}
\newcommand{\MyBibItem}[4]{%
  \refstepcounter{myref}%
  \noindent[\themyref]\ #1.\ \emph{#2}. #3. #4\par
}
\makeatother

%--------- Establecer cuatro niveles de índice --------%
    \setcounter{tocdepth}{4}   
    \setcounter{secnumdepth}{4}% numera hasta \paragraph

%%%%%%%%%%%%%%%%%%%%%%%%%%%%%%%%

\makeatletter

    %--------- Interlineado Global --------%
    \edef\GlobalBaseLineStretch{\baselinestretch}
    
    %--------- Espaciado de seccinones --------%
    \renewcommand\section{\@startsection{section}
      {1}{\z@}
      {2.5ex \@plus 1ex \@minus .2ex} % espacio antes
      {1.5ex \@plus .2ex}             % espacio después
      {\normalfont\Large\bfseries}    % formato del título
    }
    \renewcommand\subsection{\@startsection{subsection}{2}{\z@}
      {3ex \@plus 0.8ex \@minus .2ex}
      {1.5ex \@plus .2ex}
      {\normalfont\large\bfseries}
    }
    \renewcommand\subsubsection{\@startsection{subsubsection}{3}{\z@}
      {3ex \@plus 0.5ex \@minus .2ex}
      {1ex \@plus .2ex}
      {\normalfont\normalsize\bfseries}
    }
    \renewcommand\paragraph{\@startsection{paragraph}{4}{\z@}%
    {-2ex \@plus -0.5ex \@minus -.2ex}% espacio antes
    {0.8ex \@plus .2ex}% espacio después
    {\normalfont\normalsize\itshape}}% ← cursiva en lugar de negrita

    %--------- Ecuaciones --------%   
    % Variables globales para almacenar títulos y notas
    \gdef\leq@current@section{}%
    \gdef\leq@current@subsection{}%
    \gdef\leq@last@printed@section{}%
    \gdef\leq@last@printed@subsection{}%
    
    % Comando para añadir nota (invisible en el cuerpo)
    \newcommand{\eqnote}[1]{%
      \addcontentsline{leq}{leqnote}{#1}%
    }
    
    % Comando para ecuaciones
    \newcommand{\eqblock}[2]{%
      % Verificar si necesitamos imprimir el título de sección
      \ifx\leq@current@section\leq@last@printed@section\else
        \addcontentsline{leq}{leqsection}{\leq@current@section}%
        \global\let\leq@last@printed@section\leq@current@section
        \gdef\leq@last@printed@subsection{}% Reset subsección al cambiar sección
      \fi
      % Verificar si necesitamos imprimir el título de subsección
      \ifx\leq@current@subsection\leq@last@printed@subsection\else
        \ifx\leq@current@subsection\@empty\else
          \addcontentsline{leq}{leqsubsection}{\leq@current@subsection}%
          \global\let\leq@last@printed@subsection\leq@current@subsection
        \fi
      \fi
      % Ecuación
      \refstepcounter{eqnum}%
      \begin{equation*}
        #1\tag{E.\theeqnum}
      \end{equation*}%
      \addcontentsline{leq}{leqt}{[E.\theeqnum] -- #2}%
      \addcontentsline{leq}{leqm}{#1}%
    }
    
    %%% ----- Índice de ecuaciones ----------%%%
    % Tamaño para []
    \newcommand{\leqIndexFont}{\normalsize}
    \newcommand{\leqTitleFont}{\tiny}
    \newcommand{\leqNoteFont}{\tiny}
    % Cabecera de sección 
    \newcommand*\l@leqsection[2]{%
      \addvspace{25pt}% Mayor separación antes de nueva sección
      \noindent\leqIndexFont
      \makebox[\linewidth]{\rule{.38\linewidth}{0.4pt}\ \ \textbf{  \MakeUppercase{#1}}\ \ \rule{.38\linewidth}{0.4pt}}%
      \par\vspace{-10pt}
    }
    % Cabecera de subsección 
    \newcommand*\l@leqsubsection[2]{%
      \par\addvspace{15pt}%
      \noindent\leqIndexFont #1
      \par\vspace{-7pt}
    }
    % Título de ecuación 
    \newcommand*\l@leqt[2]{%
    \par\addvspace{10pt}%
      \par\noindent\leqTitleFont #1\par
    }
    % Miniatura de ecuación
    \newcommand*\l@leqm[2]{%
      \vspace{0.3\baselineskip}%
      \par\scriptsize\noindent\hspace*{1.5em}\(#1\)\par%
    }
    % Nota de ecuación 
    \newcommand*\l@leqnote[2]{%
      \vspace{0.2\baselineskip}%
      \par\noindent\leqNoteFont\hspace*{1.5em}\textbf{Nota.--} #1\par\vspace{0.5\baselineskip}%
    }
    % Generación del índice
    \newcommand{\mylistofequations}{%
      \clearpage
      \phantomsection
      % Título visual de la lista (ajústalo a tu gusto):
      {\centering\bfseries\Large Índice de ecuaciones\par}
      % Entrada en nuestro índice 'stc':
      \addcontentsline{stc}{section}{Índice de ecuaciones}%
      % Llamada a tu lista real (usa el nombre exacto de tu macro):
      \@starttoc{leq}%
    }
    
    %--------- Captura de títulos de sección --------%
    \let\leq@original@section\section
    \let\leq@original@subsection\subsection
    % Flag para saber si estamos en una sección asterisco
    \newif\ifLeqStarSection
    \LeqStarSectionfalse
    
    % Redefinir \section CON CONTROL DE ASTERISCO
    \renewcommand{\section}{%
      \@ifstar{\leq@section@star}{\@dblarg\leq@section@nostar}%
    }
    \def\leq@section@star#1{%
      \global\LeqStarSectiontrue
      \gdef\leq@current@section{#1}%
      \gdef\leq@current@subsection{}%
      \leq@original@section*{#1}%
      \global\LeqStarSectionfalse
    }
    \def\leq@section@nostar[#1]#2{%
      \global\LeqStarSectionfalse
      \gdef\leq@current@section{#2}%
      \gdef\leq@current@subsection{}%
      \leq@original@section[#1]{#2}%
    }
    % Redefinir \subsection
    \renewcommand{\subsection}{%
      \@ifstar{\leq@subsection@star}{\@dblarg\leq@subsection@nostar}%
    }
    \def\leq@subsection@star#1{%
      \global\LeqStarSectiontrue
      \gdef\leq@current@subsection{#1}%
      \leq@original@subsection*{#1}%
      \global\LeqStarSectionfalse
    }
    \def\leq@subsection@nostar[#1]#2{%
      \global\LeqStarSectionfalse
      \gdef\leq@current@subsection{#2}%
      \leq@original@subsection[#1]{#2}%
    }

    % ----- Restauración de valores Globales de estilo ----------%
    % Flags
    \newif\ifInFootnote
    \InFootnotefalse
    \pretocmd{\@footnotetext}{\InFootnotetrue}{}{}
    \apptocmd{\@footnotetext}{\InFootnotefalse}{}{}
    % Profundidad de listas (soporta anidadas)
    \newcount\MyListDepth \MyListDepth=0
    \AtBeginEnvironment{la}{\advance\MyListDepth by 1}
    \AfterEndEnvironment{la}{\advance\MyListDepth by -1}
    \AtBeginEnvironment{lnum}{\advance\MyListDepth by 1}
    \AfterEndEnvironment{lnum}{\advance\MyListDepth by -1}
    \AtBeginEnvironment{itemize}{\advance\MyListDepth by 1}
    \AfterEndEnvironment{itemize}{\advance\MyListDepth by -1}
    \AtBeginEnvironment{enumerate}{\advance\MyListDepth by 1}
    \AfterEndEnvironment{enumerate}{\advance\MyListDepth by -1}
    % Restauración diferida: hazlo al INICIO del próximo párrafo real
    \newif\if@pendingRestore
    \@pendingRestorefalse
    \newcommand{\RequestRestore}{\global\@pendingRestoretrue}
    \everypar{%
      \if@pendingRestore
        \global\@pendingRestorefalse
        \ifInFootnote\else
          \ifnum\MyListDepth=0
            \setlength{\parskip}{\GlobalParskip}%
            \setstretch{\GlobalBaseLineStretch}%
          \fi
        \fi
      \fi
    }
    % Al final de bloques:
    \AfterEndEnvironment{equation}{\RequestRestore}
    \AfterEndEnvironment{equation*}{\RequestRestore}
    \AfterEndEnvironment{la}{\RequestRestore}
    \AfterEndEnvironment{lnum}{\RequestRestore}
    % Al final de macros:
    \apptocmd{\eqblock}{\RequestRestore}{}{}
    \apptocmd{\imagenfit}{\RequestRestore}{}{}
    
\makeatother

% ===================== ToC propio con 4 niveles (único bloque) =====================
\makeatletter

% --- Escritores: añadimos una línea al .stc en cada cabecera numerada ---
% Guardamos las versiones actuales para llamar al comportamiento normal
\let\MyTOC@orig@section\section
\let\MyTOC@orig@subsection\subsection
\let\MyTOC@orig@subsubsection\subsubsection
\let\MyTOC@orig@paragraph\paragraph

% SECTION
\renewcommand{\section}{%
  \@ifstar{\MyTOC@section@star}{\@dblarg\MyTOC@section@nostar}%
}
\def\MyTOC@section@star#1{%
  \MyTOC@orig@section*{#1}% sin entrada al .stc
}
\def\MyTOC@section@nostar[#1]#2{%
  \MyTOC@orig@section[{#1}]{#2}% comportamiento normal (escribe al .toc)
  % Escribimos también nuestra línea al .stc (texto con \numberline)
  \addcontentsline{stc}{section}{\protect\numberline{\thesection}#2}%
}

% SUBSECTION
\renewcommand{\subsection}{%
  \@ifstar{\MyTOC@subsection@star}{\@dblarg\MyTOC@subsection@nostar}%
}
\def\MyTOC@subsection@star#1{%
  \MyTOC@orig@subsection*{#1}%
}
\def\MyTOC@subsection@nostar[#1]#2{%
  \MyTOC@orig@subsection[{#1}]{#2}%
  \addcontentsline{stc}{subsection}{\protect\numberline{\thesubsection}#2}%
}

% SUBSUBSECTION
\renewcommand{\subsubsection}{%
  \@ifstar{\MyTOC@subsubsection@star}{\@dblarg\MyTOC@subsubsection@nostar}%
}
\def\MyTOC@subsubsection@star#1{%
  \MyTOC@orig@subsubsection*{#1}%
}
\def\MyTOC@subsubsection@nostar[#1]#2{%
  \MyTOC@orig@subsubsection[{#1}]{#2}%
  \addcontentsline{stc}{subsubsection}{\protect\numberline{\thesubsubsection}#2}%
}

% PARAGRAPH
\renewcommand{\paragraph}{%
  \@ifstar{\MyTOC@paragraph@star}{\@dblarg\MyTOC@paragraph@nostar}%
}
\def\MyTOC@paragraph@star#1{%
  \MyTOC@orig@paragraph*{#1}%
}
\def\MyTOC@paragraph@nostar[#1]#2{%
  \MyTOC@orig@paragraph[{#1}]{#2}%
  % Si no numeras los \paragraph, cambia \theparagraph por texto plano sin \numberline
  \addcontentsline{stc}{paragraph}{\protect\numberline{\theparagraph}#2}%
}

% --- Impresor del índice propio (lee el .stc) ---
\newcommand{\MyTOC}{%
  \par\bigskip
  {\centering\bfseries\Large Índice de contenido\par}%
  \medskip
  \begingroup
    \parindent=0pt\parskip=2pt
    \@starttoc{stc}%
  \endgroup
}

\makeatother
% ===================== fin ToC propio =====================


%%%%%%%%%%%%%%


% --- Datos del tema ---
\title{Tema 3.A.25 -- La teoría neoclásica del mercado de trabajo\\
\large Análisis inter-temporal de la oferta de trabajo. Teoría del capital humano. Función de ingresos de capital humano y evidencia empírica}
\author{Marcelo Ortega}
\date{Noviembre 2025}
\newcommand{\TemaCode}{3A25}

\oldtitle{Tema 3.A.18. Teoría neoclásica de la oferta y la demanda de trabajo. Análisis intertemporal de la oferta de trabajo. Teoría del capital humano}
\changerationale{Se clarifica que la problemática central del tema es entender por qué los hogares deciden ofrecer sus horas de trabajo. Se refuerza el enfoque empírico, de especial relevancia desde la década de los años 90's y premiado con el último premio Nobel de Economía.}


\begin{document}


\maketitle
\changebox

\newpage

% --- Página de resumen y modificaciones ---
\puntosclave{ %% Escribir puntos clave Apuntes CECO
    Apuntes CECO -- 38 págs. + Anexos 

    Si bien el presente tema opta por explicar de manera especialmente detallada el enfoque neoclásico
para facilitar la comprensión al opositor de primera vuelta, en el momento del examen se
recomienda un cante equilibrado en tiempos entre los tres enfoques (en especial, la demanda de
trabajo en el modelo neoclásico estático se presentará idealmente únicamente con el fin de poder
determinar el equilibrio en el mercado). El buen opositor realizará por tanto un esfuerzo de
investigación más allá del contenido que aquí se presenta, especialmente para el tercer apartado,
para el cual sólo se aportan unas pinceladas\footnote{%
    Los principales métodos econométricos para la realización de experimentos naturales son los modelos de diferencias-en-diferencias, regresiones con discontinuidades, modelos de \textit{matching}, y el uso de variables instrumentales. Se
recomienda al opositor visitar el siguiente enlace como punto de partida: \href{https://dimewiki.worldbank.org/Quasi-Experimental_Methods}{Quasi-Experimetal Methods, World Bank Group}
}}
\vspace{1cm}

\modificaciones{
    Introducir la demostración del STFEB: Manual 
    }

\newpage
\MyTOC
\newpage

\mylistofequations
\clearpage

\MyListOfFigures
\clearpage


\pagenumbering{arabic}
\setcounter{page}{1}


\section{Introducción}\label{intro}
La teoría del mercado de trabajo, o Economía Laboral, analiza la determinación del empleo y los salarios de equilibrio
que fijan los trabajadores (que ofertan su fuerza de trabajo) y las empresas (que demandan la
fuerza de trabajo). El mercado de trabajo se puede analizar tanto desde el punto de vista:
\begin{la}
    \item Microeconómico, estudiando el comportamiento de los agentes individuales y la fijación del
equilibrio de mercado; como,
    \item Macroeconómico, estudiando las fluctuaciones cíclicas
y tendencia del mercado de trabajo así como su influencia sobre el conjunto de la economía.
\end{la}

Los modelos empleados en el este tema tienen un perfil más cercano a la Microeconomía, mientras que los estudiados en [\authorfont{Ver Temas 3.A.26, 3.A.27 y 3.A.28}] suelen asociarse a la Macroeconomía. Todo ello, con salvedades, y siendo conscientes de que no existe tal cosa corno una frontera nítida entre Macroeconomía y Microeconomía.


\subsection{Contextualización}\label{context}

\subsubsection{Relevancia}\label{importance}
El estudio del mercado de trabajo es
fundamental para tratar de comprender los mecanismos de fijación del salario y el empleo de
equilibrio, así como la transmisión de los posibles desequilibrios que genere.

\subsubsection{Historia}\label{history}
Debido a la gran importancia de la economía laboral, son muchos los economistas que se han centrado en su estudio.
\begin{lnum}
    \item Los autores neoclásicos aplicaron el principio de marginalidad y, en concreto, la teoría de la productividad marginal del trabajo para el estudio de la demanda de trabajo. [\authorfont{Ver Tema 3.A.25}]
    \begin{la}
        \item MARSHALL se limita a señalar algunos aspectos de la oferta y la demanda de trabajo que dejan ver que el mercado de trabajo se puede analizar como cualquier otro mercado – equilibrio parcial – tratando al mercado de trabajo separadamente. Aun así, no se encuentra una investigación sistemática en sus \textit{Principios de Economía} (1890).\footnote{%
            Para MARSHALL el desempleo no parecía ser un problema abrumador, lo que podría ser explicado por el contexto de su obra. Según MATTHEWS (1990):\\
“\textit{Unemployment, particularly in combination with inflation has made the functioning of the labor market a central topic in present-day economics. Unemployment has been judged as both intellectually anomalous and a social challenge. This emphasis is absent in Marshall. The social problem that disturbed his conscience was poverty; and poverty might have a number of causes, of which unemployment was only one.}”
        }
        \item HICKS, en su obra “\textit{Theory of Wages}” (1932), estudió la demanda de trabajo mediante la productividad marginal del trabajador. HICKS señala que la teoría económica no era capaz de explicar el desempleo, un fenómeno innegable de la realidad\footnote{HICKS escribe en el contexto de la Gran Depresión.}. Señala que de acuerdo con la teoría neoclásica los salarios deben reducirse en presencia de desempleo involuntario y llega a la conclusión de que tanto los sindicatos como el desempleo friccional juegan un rol importante\footnote{%
            HICKS en el marco del modelo IS-LL atribuye el desempleo a la rigidez salarial y deja entrever que tiene que ver con la forma de negociación salarial y el papel de los sindicatos. Esta rigidez salarial provoca un salario que no vacía el mercado de trabajo y por tanto genera paro involuntario.
        }. Sin embargo, la obra de HICKS fue eclipsada por la publicación de la Teoría General del Empleo, el Interés y el Dinero de KEYNES (1936).
    \end{la}
    \item Debido a la imposibilidad de explicar el desempleo\footnote{%
        Definición de según la Organización Internacional del Trabajo: “Las personas en desocupación, o personas desocupadas, se definen como todas aquellas personas en edad de trabajar que no estaban ocupadas, que habían llevado a cabo actividades de búsqueda de un puesto de trabajo durante un período reciente especificado, y que estaban actualmente disponibles para ocupar un puesto de trabajo en caso de que existiera la oportunidad de hacerlo.”
    } de estas teorías, se necesitaban nuevas contribuciones. Estas contribuciones realizadas entre la década de 1970 y la década de 1990 cambian los supuestos tecnológicos marshallianos y obtienen como resultado la emergencia de desempleo al permitir un ajuste no perfecto del empleo ante perturbaciones. [\authorfont{Ver Tema 3.A.26}]
    \item Finalmente, también era necesario ahondar en el proceso de determinación de los salarios. De manera contemporánea a los modelos anteriores, surgen modelos que explican las rigideces salariales [\authorfont{Ver Tema 3.A.27}]
\end{lnum}

\subsubsection{Problemática}\label{problem}
En el presente Tema intentaremos dar respuesta a las siguientes cuestiones:
\begin{lnum}
    \item \textit{¿Por qué los individuos deciden ofrecer su tiempo y esfuerzo en forma de trabajo?}
    \item \textit{¿Qué determinantes condicionan esta oferta de trabajo y por tanto el equilibrio en el
mercado laboral?}
\end{lnum}

\subsection{Estructura de la exposición}\label{structure}
Es razonable estructurar la exposición desde tres aproximaciones económicas distintas
pero que han de entenderse y aplicarse de manera complementaria. 
\begin{lnum}
    \item En primer lugar, el modelo
baseline es el de un enfoque neoclásico, tanto en un marco estático (modelo ocio-consumo) como
en un marco dinámico (germen de los modelos macroeconómicos del ciclo real);
\item En segundo
lugar, la teoría del capital humano crea un enfoque más específico, centrando la atención en la
motivación y las consecuencias de las tomas de decisión de trabajo y formación de los individuos; y,
\item En tercer lugar, desde los años 90, han proliferado las aproximaciones empíricas que buscan
entender las relaciones de causalidad entre las distintas políticas económicas y cambios en el
contexto por un lado, y la oferta de trabajo por otro.
\end{lnum} 

\clearpage


\section{Modelo neoclásico}
\puntosclave{
La oferta de trabajo se deriva del problema de maximización de la utilidad del consumidor
sujeto a las restricciones presupuestarias.

Las propiedades de la oferta de trabajo individual provienen de la combinación del efecto
sustitución y del efecto renta, que dan lugar una relación no monotónica entre los salarios
y la oferta individual de trabajo.

El problema de la empresa se limite a determinar el volumen óptimo de empleo para
maximizar sus beneficios.

La introducción del tiempo en un modelo dinámico no altera en esencia lo anterior, si bien
lo complementa y permite la construcción de modelos macroeconómicos microfundamentados.
}
La Teoría Neoclásica analiza el mercado de trabajo siguiendo el mismo enfoque que el caso de
un mercado donde se intercambia cualquier otro bien, es decir, derivando por tanto su función de
oferta y de demanda, que se corresponderían con el comportamiento óptimo de los agentes que
participan en el mercado.

En este sentido, se considera que:
\begin{lnum}
    \item El mercado de trabajo es perfectamente competitivo;
    \item Información transparente entre sus agentes, donde los trabajadores y las empresas son informados
perfectamente sobre los salarios y la cantidad de trabajo disponible por otras empresas;
    \item Homogeneidad del trabajo; y,
    \item No hay fricciones ni costes relacionados con el intercambio y matching de trabajadores y puestos vacantes disponibles.\footnote{%
        No hay tiempo de búsqueda de trabajo, no hay costes de desplazamiento a una entrevista de trabajo. Estas dos hipótesis de partida del modelo clásico, información perfecta y
ausencia de fricciones en el mercado de trabajo son bastante exigentes y raramente existen en la
práctica en el mercado de trabajo. De hecho, los oferentes de empleo suelen dedicar esfuerzos
significativos para encontrar un trabajo satisfactorio y la información suele ser imperfecta.
    }
\end{lnum} 

Como resulta evidentes, los supuestos del modelo neoclásicos no se corresponden con las características
que determinan el funcionamiento del mercado de trabajo.\footnote{%
    Por ello, los enfoque de la
Nueva Economía Keynesiana, intentan superar los supuestos neoclásicos y explicar el mercado
de trabajo introduciendo fricciones de distinto tipo con respecto al modelo neoclásico.
} Sin embargo, el modelo neoclásico resulta útil puesto que aporta una visión del funcionamiento de
forma sencilla y simplificada, que sirve de referencia para el análisis del equilibrio del mercado
de trabajo.

\subsection{Enfoque estático}
\subsubsection{Obtención de la oferta}

El punto de partida del enfoque neoclásico del mercado de trabajo
se basa en suponer que el individuo tiene una cantidad limitada de tiempo disponible y tiene que
elegir entre dedicarlo a un trabajo remunerado o a ocio.\footnote{%
    El salario, además, constituye un factor
fundamental en la elección. Por tanto, la oferta de trabajo depende del intercambio con la familia
dado que es una parte fundamental de la producción del hogar.
}

\paragraph{El problema del consumidor: Elección entre consumo y ocio}
La modelización de la oferta de trabajo en el marco neoclásico comienza con las decisiones tomadas por un individuo quien tiene a su disposición un conjunto de elección formado por diferentes vectores, cada uno de los cuales se compone de dos elementos:
\begin{la}
    \item El consumo de una determinada cantidad de bienes; y,
    \item El disfrute de una cantidad determinada de horas en forma de ocio.  
\end{la} 

Estas preferencias del agente vendrán representadas por medio de una función de utilidad neoclásica de buen comportamiento, $U(C,L)$\footnote{%
    La función de utilidad representa las preferencias de un individuo con respecto al ocio y al consumo. Cumple las
siguientes propiedades: las preferencias son racionales (completas y transitivas), convexas y continuas. Por ello, la
función de utilidad que las representa será monótona, continua y cuasicóncava. El individuo querrá consumir la mayor
cantidad posible de bienes y de ocio, lo que provoca que la función de utilidad sea creciente en ambos argumentos.
Como las preferencias son convexas, el individuo obtiene la misma utilidad de consumir: mucho ocio y pocos bienes o
muchos bienes y poco ocio.
}, por lo que su objetivo será elegir la cesta que le permita maximizar su función de utilidad. Además, cabe destacar que dicha combinación $C^*,L^*$ se ve también influenciada por ingresos que puedan provenir de otras fuentes ajenas al tiempo dedicado al trabajo, por tanto, distinguiremos y tendremos en cuentalo siguiente:
\begin{lnum}
    \item \textit{Fuentes no salariales}. Ingresos percibidos ajenos al mercado de trabajo, como subvenciones, donaciones, herencias, ingresos de inversiones, etc. 
    \item \textit{Ingresos salariales}. Ingresos que percibe el trabajador gracias a participar en el mercado de trabajo, dedicando unas horas al trabajo a cambio de una remuneración determinada.\footnote{%
        Este segundo punto es la esencia del problema de decisión del consumidor: tanto los bienes como el ocio le aportan utilidad. Querría tener tantos bienes y tanto ocio como fuese posible, pero para poder remunerarlos necesita trabajar, lo que supone dedicar horas de su tiempo disponible al trabajo y no al ocio. La valoración de dicho trade-off será la que determine la decisión de cuánto trabajar por parte del agente.
    }
\end{lnum}

Al igual que ocurría en la modelización neoclásica del consumidor, el agente enfrenta tanto un \textit{trade-off} como una restricción que limita sus posibilidades de elección. De esta forma, su problema de optimización puede representarse de la siguiente manera:
\eqblock{
\begin{aligned}
&\max_{\{C,L\}}{U(C,L)}\qquad \text{s.a.} \quad \begin{cases}
    C &\leq w(L_0-L)\;+\; R\\
    L_0&\ge L
\end{cases}
\hspace{2em};\text{Donde: }\quad
\begin{cases}
    L_0&\equiv\quad \text{\scriptsize Dotación de tiempo}\\
    L&\equiv\quad \text{\scriptsize Horas de ocio}\\
    w&\equiv\quad \text{\scriptsize Salario real por hora, $w_h=w(L_0-L)\equiv$ Renta salarial total}\\
    R&\equiv\quad \text{\scriptsize Renta que el agente obtiene fuera del mercado de trabajo}
\end{cases}\\[1.5em]
&\Longrightarrow \text{Por tanto,}\hspace{6em}
\boxed{\begin{aligned}
    &\max_{\{C,L\}}{U(C,L)}\\[1.5em]
    &\text{s.a.}\quad C\;+\;wL\leq \underbrace{R_0}_{\text{\scriptsize Renta potencial, $wL_0\;+\;R$}}
\end{aligned}}
\end{aligned}
}{Oferta neoclásica. Problema del consumidor: Maximización Consumo-Ocio}

Como vemos, a la tradicional restricción presupuestaria se le añade una restricción adicional, que es la restricción temporal, ya que el agente dispone de una cantidad determinada de horas, y no puede dedicar más que las que tiene, por lo que tendrá que decidir cómo repartirlas, porque no puede aumentarlas. Calculando las CPO's:
\eqblock{
\begin{aligned}
    \text{De,} \quad &\mathcal{L}=U(C,L)+\mu[C+wL-R_0]\\[1em]
    &\text{Obtenemos,}\\[1.5em]
    &\begin{cases}
    &\text{(1)}\qquad \frac{U_L(C^*, L^*)}{U_C(C^*, L^*)}=\frac{\partial U/\partial L}{\partial U/\partial C}|_{(C^*,L^*)}=w\\[1em]
    &\text{(2)}\qquad C^*+w\;L^*=R_0
    \end{cases}
\end{aligned}
}{Oferta Neoclásica. Problema del consumidor: CPO's}
 
Es decir, en el óptimo el agente se ubicará en aquel punto en el que la RMS entre ocio y consumo, $RMS_L^C$, o valoración subjetiva, se iguala al salario real o valoración objetiva.\footnote{Si no se ubicase en este punto sería posible aumentar su utilidad.}  Gráficamente, lo que busca el agente es situarse en la Curva de Indiferencia más alejada posible del origen, cumpliendo con su restricción presupuestaria:
\imagenfit{SolMaxS_L.png}{$E$: Óptimo en la oferta individual de trabajo}. Restricción y Curvas de Indiferencia{Apuntes ICEX-CECO. Tema 3.A.25}

Así, obtendríamos la Función de oferta de trabajo correspondiente como la diferencia entre la dotación de tiempo y la demanda de ocio, $h^*=L_0-L^*$.\footnote{%
    Esta función de oferta de trabajo nos indicará, por tanto, cuál es la cantidad de horas de trabajo que el agente ofrece en el mercado para un salario real y unas renta potencial determinados
} Se le denomina habitualmente como oferta de trabajo Marshalliana o "no compensada":
\eqblock{
\begin{aligned}
&\text{\textbf{Definición:} Función de oferta de trabajo Marshalliana o no compensada}\\[2em]
&\hspace{10em}\boxed{
L^*=\Gamma(w,R_0)}
\end{aligned}
}{Oferta Neoclásica. Función de oferta Marshalliana}

Una de las cuestiones que mayor interés ha suscitado en la literatura es la forma que adopta esta función de oferta de trabajo Marshalliana o no compensada, y va a depender de los efectos sustitución y efecto renta que se producen al variar el salario.

\paragraph{Variaciones en la Función de Oferta Marshalliana}
Démonos cuenta que depende de dos variables: salario real, $w$, y renta potencial, $R_0$. Esto da lugar
a una relación no monótona entre salario y oferta de trabajo.

\textsc{Cambios en la renta no salarial. El Salario de reserva}\\
El hecho de que la curva de indiferencia sea convexa implica que la relación marginal de
sustitución entre consumo y ocio sea
decreciente conforme nos movemos hacia la izquierda y hacia abajo en la curva de indiferencia.
Por tanto, el individuo sólo ofrecerá horas de trabajo cuando se cumpla que la utilidad marginal
del trabajo con respecto a la utilidad marginal del ocio sea mayor que este salario $(U_L /U_C > w)$.
La relación marginal de sustitución en este punto se llama salario de reserva, y se define como:
$$w_A=\frac{U_L(R, L_0)}{U_C(R,L_0)}$$
Depende de la forma de la función de utilidad y del nivel del ingreso no salarial $R$. El salario de
reserva determina la condición de participación en el mercado de trabajo, de forma que, si el
salario cae por debajo del salario de reserva, el trabajador no participará en el mercado de trabajo.En la medida en que la renta no salarial. $R$, aumente, siendo el ocio un bien normal, el salario de
reserva también se incrementará, manteniéndose la utilidad constante. Por ello, se suele decir que
la renta no salarial tiene un efecto desincentivador en la decisión de participación en el mercado
de trabajo.
\imagenfit{SolMaxS_L.png}{$E\to E'$: Incremento en la renta no salarial, $R\to R_1$}. Restricción y Curvas de Indiferencia{Apuntes ICEX-CECO. Tema 3.A.25}

\textsc{Cambios en el salario real}\\
Como hemos visto, el ocio sería un bien normal cuando su demanda se incrementa al incrementarse la renta no salarial, y sería un
bien inferior si su demanda decrece al incrementarse la renta no salarial. Sin embargo, la relación entre oferta de trabajo y salario real se complica. 

Démonos cuenta que el salario real, $w$, puede considerarse un precio, pues es el coste de oportunidad del ocio, $L$. Por tanto, atendiendo a la Tª Neoclásica de la Demanda del Consumidor [\authorfont{Ver Tema 3.A.9}], un aumento de un precio tiene un efecto sustitución (derivado del cambio en los precios relativos), y un efecto renta (por el cambio en el poder adquisitivo de la renta). Sin embargo, ahora la renta del agente no es exógena, ya que se compone en parte del salario real y, por tanto, un aumento del salario real la incrementa, generando un efecto adicional que es el efecto renta-dotación. Tenemos pues tres efectos que debemos valorar ante cambios en el salrio real:
\begin{la}
    \item Un efecto sustitución, por el encarecimiento relativo; y,
    \item Dos efectos rentas:
    \begin{lnum}
        \item Efecto Renta (normal), relativo al poder adquisitivo del consumidor;
        \item Efecto Renta-Dotación, efecto endógeno, generado por la revalorización de la dotación inicial de tiempo ($w_0L_0<w_1L_0$), gracias al aumento del salario real.
    \end{lnum}
\end{la}  

El efecto sustitución tiene el signo habitual, no positivo, pues cuando aumenta el salario, el ocio se
encarece relativamente y disminuye su consumo y, por tanto, aumenta la oferta de trabajo. El
signo del efecto renta depende de si el ocio es un bien normal o inferior:
\begin{la}
    \item Si el ocio es un bien normal, cuando aumente el salario aumentará su consumo y, por tanto,
disminuirá la oferta de trabajo.
    \item Si el ocio es un bien inferior, sucedería lo contrario. Lo anterior también es aplicable tanto al
efecto renta-dotación como al efecto renta-directo.
\end{la}

Teniendo en cuenta que renta total es la suma de la renta salarial y no salarial, $R_0 = wL_0 + R$, para cuantificar el sentido de estos efectos, calculamos la derivada de la demanda de ocio en el óptimo:

\eqblock{
\begin{aligned}
    &\text{(Por la Regla de la cadena) De:}\\[2em]
    &\hspace{4em}\underbrace{L^*=\Gamma(w,R_0)}_{\text{\scriptsize Función Oferta Marshalliana}} \quad \to \quad \frac{\partial L^*}{\partial w}=\underbrace{\Gamma_1}_{\text{\scriptsize Efectos clásicos: $ES+ ER$}}+\underbrace{\Gamma_2\frac{\mathrm{d}R_0}{\mathrm{dw}}}_{\text{\scriptsize Efecto Renta-dotación}};\\[2em]
    &\text{Siendo,} \qquad \Gamma_1=\frac{\partial L^*}{\partial w},\qquad \Gamma_2=\frac{\partial L^*}{\partial  R_0}\quad \text{y}\quad  \frac{\mathrm{d}R_0}{\mathrm{d} w}=R_0'(w)=L_0 
\end{aligned}
}{Oferta neoclásica. Derivada de la Función de Oferta Marshalliana respecto al salario}

Por lo que hay varios elementos que hay que tener en cuenta:\footnote{%
    Otra forma de representar esto sería mediante la ecuación de Slutsky:
$$\frac{\partial L^*}{\partial w}=\frac{\partial (L^*)^c}{\partial w}+(L_0-L)\frac{\partial L^*}{\partial R_0}$$
Démonos cuenta que en la ecuación de SLUTSKY, el Efecto Renta global, conforme se reduce el ocio (aumenta el trabajo) va siendo mayor, por lo que es de esperar que cuando la oferta de trabajo ya es muy elevada éste acabe dominando al primero y que al ser el ocio un bien normal, la relación pase a ser positiva y, por tanto, la relación trabajo-salario pase a ser negativa.
}
\begin{lnum}
    \item $\Gamma_1=\frac{\partial L^*}{\partial w}$ será la combinación del Efecto Sustitución y del Efecto Renta indirectos. 
    \begin{la}
        \item El Efecto Sustitución será \textit{no-positivo} por definición\footnote{%
            Cuando el ocio se encarece relativamente, manteniendo la renta real constante, nunca aumentaremos su consumo
        }; mientras que,
        \item El signo del Efecto Renta dependerá de si estamos ante un bien normal o un bien inferior.
    \end{la}
    \item $\Gamma_2\;R_0'(w)$, corresponde al Efecto Renta-Dotación (efecto directo) y el signo dependerá de si el ocio es un bien normal o un bien inferior.\footnote{Evidentemente, $R_0' >0$, ya que es la dotación inicial de tiempo}
\end{lnum}

\imagenfit{Efectos-VarRentaReal.png}{Ajuste ante una variación del Salario Rela, $w\to w_1$}{Apuntes ICEX-CECO. Tema 3.A.25}

Suponiendo que el ocio es un bien normal\footnote{%
    Generalmente en la literatura se asume que el ocio es un bien normal, y se acepta el hecho de que la oferta de trabajo es positiva con el salario real para los primeros tramos, pero que llega un punto en el que su curvatura cambia.
} y para representar el efecto sustitución y el efecto renta
indirecto resulta útil realizar un análisis en dos etapas:
\begin{lnum}
    \item En la primera, se analizan el efecto
sustitución y el efecto renta indirecto para lo que se supone que a pesar de que el salario aumenta,
la renta potencial no varía (es decir, se asume que la renta no-salarial disminuye, $\downarrow R$). Esto implica
que nos movemos en la restricción presupuestaria $\overline{R_0A_1}$ , con pendiente $w_1$. Esta restricción
presupuestaria es compensada con $R_c=R - (w_1 - w)L_0$ y la pendiente en este sería $w_1$;\footnote{%
    Desde el equilibrio inicial $E$, el efecto sustitución representa la variación en el consumo de ocio
al variar el salario, manteniendo el nivel de utilidad (curva de indiferencia original), de forma que
vendría dado por el paso de $E$ a $E'$.Este desplazamiento se corresponde con el cambio
Hicksiano (o compensado) de la oferta de trabajo, obtenido de minimizar el gasto del
consumidor bajo la restricción de conseguir un nivel de utilidad dado. Por ello, el efecto
sustitución generaría una reducción de tiempo de ocio.
Sin embargo, partiendo de $E'$ y, suponiendo que el salario se mantiene en él nivel de $w_1$, el efecto
renta indirecto hace que el equilibrio del consumidor se desplace hasta $E"$. En este nuevo
equilibrio, la renta es menor y la utilidad también es menor. En términos netos, un aumento del
salario reduce el ocio y aumenta la oferta de trabajo. Por lo tanto, el efecto sustitución y el efecto
renta indirecto producen el mismo resultado: un incremento en los salarios conduce a una
disminución del tiempo asignado a ocio, lo que conlleva un incremento en la oferta de trabajo. Si
el ocio es un bien nonnal, el efecto sustitución genera un incremento en la oferta de trabajo ante
un incremento de la renta mientras que el efecto renta lo reduciría. En definitiva, del punto $E$ pasamos al punto $E''$ por el efecto conjunto del efecto sustitución y renta indirecto. En términos
netos, un aumento del salario reduce el ocio y aumenta la oferta de trabajo. El efecto sustitución
y el efecto renta indirecto producen el mismo resultado: un incremento en los salarios conduce a
una disminución del tiempo asignado a ocio, es decir, un incremento en la oferta de trabajo. Por
tanto, tendremos que $\Gamma_1<0$ si el ocio es un bien normal.
} y,
\item En
la segunda etapa, se analiza el efecto del aumento de la renta potencial desde $R_0$ a $R_1$ ya que al
aumentar el salario el valor de la dotación total de tiempo del agente también se incrementa.\footnote{%
    El aumento del salario altera la
restricción presupuestaria del individuo aumentando la renta total (revalorizando la dotación de
tiempo). La restricción presupuestaria final, tras el aumento del salario sería $\overline{R_1A}$. resultado del
desplazamiento paralelo de $\overline{R_0A_1}$ ,, con una pendiente mayor que la restricción original al
incrementarse el salario, esto es, pasar de $w\to w_1$. Hemos supuesto que el ocio es un bien normal,
luego el aumento de la renta total ($R_0 \to R_1$) provocaría un mayor consumo de ocio, por
lo que pasaríamos del punto $E''\to E_1$.
El efecto renta directo se representa por la derivada parcial $\Gamma_2$. Si el ocio es un bien normal,
entonces por definición es positiva y cualquier incremento en los salarios conduce a un
incremento del consumo de ocio y a un descenso en la oferta de trabajo.
}
\end{lnum} 

Como consecuencia de este análisis es posible concluir que el efecto neto de un aumento del
salario depende, por tanto, de si el efecto renta-dotación compensa o no al efecto sustitución y al
efecto renta-indirecto, que operan en sentido contrario al primero. La siguiente figura representaría la función de oferta de trabajo. En función de si el ocio es un
bien normal o inferior, y de la magnitud de los efectos renta y sustitución, su pendiente será
positiva o negativa. Una vez que el salario supera el salario de reserva ($w_A$), el individuo participa
en el mercado de trabajo. Siempre que el ocio sea un bien normal, y el efecto sustitución sea
mayor en valor absoluto que el efecto renta global, la curva de oferta tendrá pendiente positiva.
Pero el efecto renta total aumenta con el salario y es razonable considerar que cuando el salario
alcanza un determinado nivel, entonces domina al efecto sustitución.\footnote{%
    El aumento de la oferta de
trabajo ante aumentos del salario dados es cada vez menor, hasta que la pendiente comenzaría a
disminuir, teniendo la curva de oferta de trabajo pendiente negativa.
}

\imagenfit{VarSalarioReal-Salarioreserva.png}{Curva de Engel -- Oferta de trabajo, $L-L_0$, respecto al salario real, $w$}{Apuntes ICEX-CECO. Tema 3.A.25}

\begin{specialblock}{Decisiones intrafamiliares}
    Las decisiones sobre el mercado de trabajo se toman, a menudo, influidas enormemente por el ámbito familiar.\footnote{%
        Los modelos de elección intrafamiliar suelen ofrecer una explicación más amplia de las decisiones de participación en el mercado de trabajo, dado que en el modelo neoclásico la decisión de participación depende del salario de mercado respecto al de reserva para cada individuo. Sin embargo, las decisiones de sus miembros suelen estar conectadas, y la variación en la renta de un individuo puede no sólo afectar a su propia renta sino también a la de los otros miembros de la familia. De hecho, lo más habitual es que la caída en el salario de un miembro provoque que haya trabajadores adicionales que entren en el mercado de trabajo. Desde el punto de vista empírico, estos efectos se han observado, aunque con limitados resultados.
    } El análisis de las decisiones familiares se ha venido estructurando en tomo a dos áreas principales:
    \begin{la}
        \item La perspectiva unitaria, que considera a la familia compuesta por dos miembros cuyas preferencias están presentes por medio de la función de utilidad común, $U(C, L_1, L_2)$, donde $L_1$ sería el tiempo que tiene uno de los miembros, y $L_2$ el otro. En esta función de utilidad, la distribución de las ganancias no es relevante, es decir, no es preciso conocer qué miembro de la familia recibe más renta. Sin embargo, éste es un supuesto que genera mucho debate.
        \item El modelo colectivo, que considera que la elección del hogar debe ser prioritaria a la elección individual. \authorfont{CHIAPPORI} (1992) desarrolló este modelo de decisión en el cual el hogar permite deducir el consumo individual, usando la oferta de trabajo individual y el consumo total del hogar, variables ambas que se pueden cuantificar.
    \end{la} 
\end{specialblock}

\paragraph{Obteción de la curva de Oferta Agregada}

Finalmente, podemos obtener la oferta agregada de mercado llevando a cabo una suma horizontal de las ofertas individuales. Cabe destacar que,
\begin{lnum}
    \item La oferta de trabajo comienza en un nivel salarial superior a cero: salario de reserva ($w_A$).
    \item Importancia de la elasticidad de la oferta de trabajo. La forma concreta que adopte la oferta de trabajo va a ser fundamental a la hora de realizar recomendaciones de política económica. Por ello cobra una especial importancia el concepto de elasticidad salario de la oferta de trabajo, que nos indica cuál es la variación porcentual de las horas de trabajo ofertadas ante un aumento de un 1\% del salario real.\footnote{%
        Así, una elasticidad positiva nos indicaría que la oferta agregada, al menos en el tramo analizado, tiene pendiente positiva. Pero no solo es importante para conocer la pendiente, sino sobre todo es importante conocer su cuantificación, pues como veremos esto tiene importantísimas implicaciones de política económica. Por ello, no es de extrañar que a nivel empírico se haya prestado una especial atención a esta cuestión, con numerosas estimaciones al respecto, siendo posiblemente el área de la economía laboral que más atención empírica ha recibido en los últimos 30 años.\\ Un survey realizado por BLUNDELL y MaCURDY destaca que la mayor parte de trabajos en la literatura encuentran una elasticidad de la oferta de trabajo muy reducida, prácticamente nula. Aunque se señala que para los hombres parece ser ligeramente negativa, y para las mujeres ligeramente positiva, la conclusión a destacar es que ésta es prácticamente igual a cero (-0,1 para hombres y 0,2 para mujeres).
    }
\end{lnum} 

\subsubsection{Obtención de la curva de demanda}
La teoría de la demanda de trabajo está profundamente relacionada con la producción de factores.
Sin embargo, no sólo depende del coste del trabajo, sino también del coste de otros factores de
producción y otros factores que determinan el beneficio del empresario.\footnote{%
    El otro lado del mercado está formado por la demanda de trabajo que obtendremos mediante la resolución del Problema del Productor [\authofont{Ver Tema 3.A.11}].
} Estos otros factores
pueden ser la eficiencia del desempeño de los trabajadores y el precio al que se venden los bienes\footnote{%
    La eficiencia del trabajo depende de la tecnología disponible y las cantidades de otros factores de producción, como
capital o energía, utilizados por las empresas y de las características de cada trabajador.
} La empresa contratará trabajadores (o aumenta el tiempo trabajado por ellos) para realizar su
actividad productiva, combinándolos con otros factores de producción, si la renta generada por el
trabajador supera el coste de contratarlo. Por tanto, la demanda de trabajo de una empresa
maximizadora de beneficios depende del salario real y del poder de mercado de la empresa.

Tal y como señala la literatura microeconómica, es útil realizar una distinción entre el corto plazo y el largo plazo.

\paragraph{Corto plazo}
Se caracteriza porque al menos uno de los factores productivos es fijo y al menos otro es variable. En este caso vamos a asumir al trabajo como factor variable y al capital como factor fijo. Es decir, las empresas ajustarán la contratación de trabajadores para maximizar su beneficio teniendo un capital determinado.

\eqblock{
\begin{aligned}
    &\max_{\{L\}}{\Pi(\bar K, L)}=F(\bar K, L) - wL-r\bar K  \quad \text{s.a.} \quad \Pi\ge0\\[2em]
    &\text{Obtendremos la CPO:}\\
    &PMg_L=w \qquad \equiv \qquad PMg_L=CMg_L;\\[2em]
    &\text{Y, finalmente, }\\
    &\underbrace{L^D=\Gamma{(w, \bar K)}}_{\text{\scriptsize Función de Demanda Marshalliana de trabajo}}\qquad \text{Donde, }\quad 
        \partial L^D/\partial w<0
\end{aligned}
}{Demanda neoclásica. Problema del Productor}

Vemos entonces cómo ésta tiene pendiente negativa, puesto que la PMG es decreciente, por lo que una reducción del salario real implica una reducción de la productividad marginal y para ello hay que contratar más trabajadores. También vemos que se desplazará horizontalmente ante cambios en el stock de capital existente. 
 
\paragraph{Largo plazo}
Cuando consideramos el largo plazo conviene aclarar que el factor capital se convierte también
en flexible. Por ello, deberemos analizar los efectos sustitución y los efectos escala,
que permitirán deducir las propiedades de las demandas de factores condicionadas y no
condicionadas:
\begin{lnum}
    \item En una primera etapa, se analiza el efecto sustitución que determina cómo varía
la combinación óptima de factores, para alcanzar un nivel de producción dado, cuando cambian
los precios relativos de los factores, lo que nos permite obtener la función de demanda
condicionada de factor trabajo; y,
    \item En una segunda etapa, se estudia el efecto escala, que mide cómo
varía la demanda de factores, ante cambios en el volumen de producción, manteniendo la
proporción en la que ambos se combinan.\footnote{%
    La función de producción se caracteriza por la existencia de productividades marginales de cada uno de los factores positivas pero decrecientes, es decir $F'(L)>0, \;F'(K)>0\; y \; F''(L)<0\;, \; F''(K)<0$
}
\end{lnum}  

\textsc{Primera Etapa: Minimización del coste}\\
El problema de la empresa en el largo plazo será buscar, por tanto, la combinación óptima de
factores de producción que minimice el coste de producción de la empresa sujeto a alcanzar un determinado nivel de producción $Y$
\eqblock{
\begin{aligned}
    &\min_{\{K,L\}} wL-rK  \qquad \text{s.a.} \quad F(K,L)\ge Y\\[2em]
    &\text{Obtendremos de las CPO's:}\\[0.5em]
    &\begin{cases}
    &\text{(1)}\quad PMg_L=w \qquad \equiv \qquad PMg_L=CMg_L;\\
    &\text{(2)}\quad PMg_K=r \qquad \equiv \qquad PMg_K=CMg_K;\\[1em]
    &\text{(3)}\quad \frac{PMg_L}{w}=\frac{PMg_K}{r} \qquad \Longrightarrow \qquad \frac{PMg_L}{PMg_K}=\frac{w}{r}=(RMST)_K^L\\
    \end{cases} \qquad \Longrightarrow \qquad \underbrace{\bar L(Y, \frac{w}{r})\quad \text{y}\quad \bar K(Y, \frac{w}{r})}_\text{\scriptsize Demandas condicionadas de factores productivos}\\[2em]
    &\text{Y, finalmente, }\\
    &\Longrightarrow\text{Función de costes de la empresa:}\quad C(w, r, Y)
\end{aligned}
}{Demanda neoclásica. Primera etapa: Minimización del Coste}

\textsc{Segunda Etapa: Maximización del beneficio}\\
Cuando no existen restricciones para la empresa sobre alcanzar un determinado nivel de
producción, la demanda de factores no está condicionada y la empresa va a elegir un nivel de
producción que maximice su beneficio.

\eqblock{
\begin{aligned}
    &\max_{\{K,L\}} \Pi(w,r,Y)=P\;Y-C(w,r,Y)  \qquad \text{s.a.} \quad F(K,L)\ge Y\\[2em]
    &\text{La CPO:} \quad \frac{\partial \Pi}{\partial Y}=0\\[0.5em]
    &\Longrightarrow\text{Óptimo de $Y$ satisface:}\quad P(Y)=v\;C(w,r,Y), \quad \text{\scriptsize donde $\quad v=1/(1+\eta_Y^P) \sim$ si en Eq. $\equiv (P=CMg):\;v=1$}\\[1.5em]
    &\text{A partir de la resolución del problema, obtendremos:}\\[0.5em]
    &\begin{cases}
        \bar L=L(\frac{w}{p}, v, K, F)=\frac{\partial F}{\partial L}=v\frac{w}{p}\\[0.5em]
        \bar K=K(\frac{r}{p}, v, L, F)=\frac{\partial F}{\partial K}=v\frac{r}{p}
    \end{cases} \qquad \text{Demandas no condicionadas (o de l/p) de $ff.pp$}
\end{aligned}
}{Demanda neoclásica. 1era etapa: Minimización del Coste}

A pesar de que esto solo constituye una orientación sobre la resolución del problema del productor a largo plazo\footnote{%
    Para más información y derivación analítica completa [\authorfont{Ver Apuntes ICEX-CECO Tema 3.A.25}]
}, sí podemos apreciar que existe una diferencia principal respecto a la demanda a corto plazo. La demanda a largo plazo del factor trabajo va a depender de su grado de sustituibilidad respecto al capital, que vamos a obtener a través de las demandas condicionadas de la oferta de trabajo (primera etapa, $\bar\eta_w^L$) y, además, de los posibles efectos de escala que pueda presentar la tecnología de producción. Por tanto, podríamos expresar la elasticidad del empleo con respecto al salario de una empresa dada como:

\eqblock{
\begin{aligned}
\underbrace{\eta_w^L}{\text{\scriptsize Elasticidad de la demanda de trabajo, $l/p$}}=\quad \underbrace{\bar\eta_w^L}_{\text{\scriptsize Efecto sustitución dado $\bar Y$}}\quad +\quad \underbrace{\bar\eta_Y^L\eta_w^Y}_{\text{\scriptsize Efecto escala}}\\[2em]
\text{Donde, }
    \begin{cases}
        \text{Efecto escala:} \quad \downarrow w\;:\quad \downarrow C(q)\;\to\;\uparrow F(K,L)\;\to\;\uparrow K^D \;\wedge\; \uparrow L^D\\[0.5em]
        \text{Efecto sustitución:} \quad\overline{F}\;\wedge\; \downarrow w\;:\quad \uparrow L^D 
    \end{cases}\\
\text{Y,}\qquad (\eta_{w}^L)^{lp}>(\eta_{w}^L)^{cp} \qquad \text{Por que:}
\begin{cases}
        \text{Corto plazo:}\quad \eta_w^L= \bar\eta_w^L\\[0.5em]
        \text{Largo plazo:} \quad \eta_w^L= \bar\eta_w^L+\bar\eta_Y^L\eta_w^Y
    \end{cases}
\end{aligned}
}{Elasticidades de la demanda de trabajo: Efecto Sustitución y Efecto Escala}

En definitiva, el efecto sustitución hace referencia a la elección de un factor sobre otro para
obtener un determinado nivel de producción, mientras que el efecto escala representaría la
capacidad para alterar el nivel de producción manteniendo las mismas proporciones de los dos
factores de producción. El efecto sustitución es siempre negativo (si $\uparrow w$),\footnote{%
    Dado
que para un nivel dado de producción, un incremento en el coste del trabajo siempre conduce a
una menor utilización de este factor.
} mientras que el efecto escala
representado es ambiguo (a priori).\footnote{%
    Excepto en el caso de una función de producción homogénea,
donde es necesariamente positivo (si $\downarrow w \to \downarrow L^D$).
}

Bajo los supuestos establecidos ambos efectos van en el mismo sentido bajo, lo que potencia la contratación para una misma reducción de salario.  Esto nos llevará por tanto a que la demanda de trabajo a largo plazo sea más elástica que la demanda a corto plazo.\footnote{%
    Será más elástica cuanto mayor sea la elasticidad de sustitución entre factores, pues más fácil será sustituir capital por trabajo, y mayor será el efecto sustitución.
}\textsuperscript{,}\footnote{%
    Desde el punto de vista empírico este tipo de trabajos basados en la teoría estática de la demanda
de trabajo suelen ir orientados a estimar las elasticidades cruzadas y de sustitución que hemos
mencionado anteriormente. El problema que conllevan este tipo de estimaciones se suele centrar
en elegir las funciones de producción y costes que se van a utilizar. La más habitual suele ser
considerar una función de producción del tipo Cobb-Douglas (1928) con dos factores de
producción o una función de elasticidad de sustitución constante (Arrow et al., 1961). Para las
funciones de costes, se suele optar por utilizar la función generalizada de Leontief (Diewert,
1971), o la función de coste traslog (Christensen, Jorgenson y Lau, 1973).
}



\paragraph{Obtención de la curva de demanda agregada}
Destacar también que la demanda agregada no es la suma horizontal de las demandas individuales. Esto se debe a que se genera una externalidad pecuniaria. Así por ejemplo, si cae el salario de los trabajadores, todas las empresas querrían contratar más trabajadores y producir más, pero esto llevaría a un aumento de la oferta del producto y una caída del precio del bien que desplazaría la demanda de trabajo. Por ello, la demanda agregada será más inelástica que la suma de las demandas individuales.\footnote{%
    Al igual que ocurría en el caso de la oferta de trabajo, un concepto de gran importancia a la hora de diseñar la política económica será el concepto de la elasticidad de la demanda de trabajo. Se define como la variación porcentual de la demanda de horas trabajadas ante una variación de un 1\% en el salario. Una obra importante en este sentido es el survey realizado por HAMERMESH, en el que se encuentra que la elasticidad de la demanda a corto plazo se ubica entre $-0,4$ y $-0,5$, y que la elasticidad a largo plazo alcanza el $-1$. Por tanto, vemos que estaríamos ante una demanda también inelástica, pero que es más elástica que la oferta de trabajo, lo que como veremos será relevante.
}

\subsubsection{Equilibrio}
En el siguiente gráfico se representa el equilibrio en $E^*$. El punto de equilibrio es el de corte de oferta y demanda de trabajo. Se trata de un equilibrio único y estable que vacía el mercado. 

La cantidad de trabajo $L^*$ equivale al pleno empleo porque el único paro que existe es el friccional.
Si existe paro, es voluntario ya que, al salario de equilibrio $(w/P)^*$ todo el que quiere trabajar
encuentra trabajo. Como el salario es perfectamente flexible, va a permitir volver al equilibrio
tras cualquier perturbación o ante cualquier desequilibrio. 

\imagenfit{Eq-NeoclasicoTrabajo.png}{Equilibrio en el Mercado de Trabajo. Modelo Neoclásico del Mercado de Trabajo}{Apuntes ICEX-CECO. Tema 3.A.25}

\paragraph{Propiedades del equilibrio}
Por lo tanto, en el equilibrio, la intersección de ambas funciones determinará la cantidad de empleo y el salario que existirá en el mercado de trabajo, de tal forma que:
\begin{lnum}
    \item Todos los consumidores maximizan su utilidad sujeto a la restricción presupuestaria y a la restricción temporal;
    \item Todas las empresas maximizan sus beneficios sujeto a la tecnología; y,
    \item Los mercados se vacían, de manera que la oferta de trabajo agregada se iguala a la demanda de trabajo agregada. 
\end{lnum}

\subsubsection{Implicaciones de política económica}
Una vez que hemos analizado la formación de este modelo neoclásico del mercado de trabajo conviene preguntarnos. ¿Qué implicaciones de política económica podríamos extraer de él? Fundamentalmente, debido a su naturaleza de mercado walrasiano en el que no existe desempleo involuntario la intervención pública al respecto se encontraría bastante limitada. 

Sin embargo, sí que podemos utilizar el modelo para analizar qué ocurriría con 2 cuestiones fundamentales que suelen afectar al mercado de trabajo:
\begin{lnum}
    \item Variaciones en el Salario Mínimo Interprofesional (SMI); y,
    \item Incidencias de un cambio fiscal, en particular, Impuestos sobre la Renta (IRPF)
\end{lnum}

Como hemos dicho a lo largo de la exposición dos conceptos fundamentales que van a afectar de forma significativa a los efectos de estas medidas sería la elasticidad de la oferta de trabajo y la elasticidad de la demanda de trabajo. Hemos visto cómo la evidencia apunta a que la elasticidad de la oferta es prácticamente nula, mientras que la demanda, aunque es bastante inelástica, es más elástica que la oferta.

\paragraph{Variaciones en el SMI}
Supongamos que nos estamos moviendo en un contexto competitivo, de forma que el mercado de trabajo se estaría vaciando a un salario de equilibrio determinado. Sin embargo, se considera que sería beneficioso para el país si se incrementase el salario mínimo por encima del salario de equilibrio. La teoría neoclásica dice que esto generaría desempleo ya que habría exceso de oferta, pero aquí es donde las elasticidades juegan un papel fundamental. 
Más concretamente, la elasticidad de la oferta de trabajo es prácticamente nula, por lo que el aumento del salario no tendría un impacto positivo sobre las horas de trabajo. Por su parte, la elasticidad de la demanda de trabajo es reducida, lo que indica que la reducción en la demanda de trabajo sería escasa. 

\textbf{Introducir gráfico alteración del SMI}
 
Ahora bien ¿cómo de escasa? A corto plazo, la evidencia apunta a que la elasticidad de la demanda se sitúa aproximadamente en el $-0,5$. Eso significa que un aumento del salario mínimo del 10\% generaría una reducción de las horas de trabajo demandadas del 5\%.

\paragraph{Variaciones del IRPF}
Otra medida que suele ser especialmente importante en el análisis de política económica es las modificaciones del IRPF para los trabajadores. Por ejemplo, supongamos el caso de un Gobierno que está considerando subir este impuesto ya que necesita incrementar su recaudación. ¿Qué encontraremos aquí? Lo que nos dice la teoría de la imposición es que los impuestos se trasladan según las elasticidades de ambos lados del mercado, teniendo un impacto mayor sobre el lado más inelástico.

\textbf{Introducir gráfico alteración del SMI}
 
En nuestro caso, hemos visto que la elasticidad de la oferta era prácticamente nula, mientras que la elasticidad de la demanda de trabajo era mayor. Por tanto, lo que nos diría la teoría económica es que la subida de este impuesto tendría un impacto casi total sobre los trabajadores, sin que hubiese apenas traslación sobre las empresas. 
Al mismo tiempo, en función de las elasticidades habrá mayor o menor destrucción de empleo. Cuanto más inelásticas sea la oferta/demanda de trabajo menos destrucción de empleo se generará, pues las variaciones en las cantidades serán ínfimas. En nuestro caso hemos visto que ambos lados del mercado se caracterizan por ser muy inelásticos, por lo que es previsible que esta figura impositiva no tenga un elevado impacto en el empleo de equilibrio. 

Por tanto, como vemos, la modelización del mercado de trabajo bajo la perspectiva neoclásica básica nos permite observar que bajo los supuestos establecidos el equilibrio alcanzado garantiza la no existencia de desempleo involuntario. Conceptos como son la elasticidad de la demanda y de la oferta son cruciales a la hora de diseñar la política económica, pues la principal bondad de este marco neoclásico es que su elevada sencillez nos permite analizar de forma muy fácil los resultados obtenidos de implementar diferentes medidas.


\subsection{Enfoque dinámico}
Una posible extensión del modelo neoclásico que, de hecho, se encuentra en la base de los
modelos de la Nueva Macroeconomía Clásica y de los Ciclos Reales\footnote{%
    LUCAS y RAPPING (1969) intentan reconciliar en su trabajo las aparente divergencias en la modelización de la oferta
de trabajo entre la teoría neoclásica y el enfoque del empleo keynesiano.
}, es el análisis las decisiones
del individuo entre consumo y ocio en un contexto intertemporal. Esto nos permite introducir los
fenómenos de sustitución intertemporal de ocio y consumo, con los que el agente responde a
shocks transitorios y permanentes de sus flujos de renta. 

Para ello, analizaremos en primer lugar
las novedades que introduce el enfoque dinámico en la oferta de trabajo respecto al análisis
estático y, posteriormente, en la demanda de trabajo. No obstante, las
conclusiones que arroja la introducción de un horizonte temporal en la determinación de la oferta
y la demanda del mercado de trabajo no modifican sustancialmente las conclusiones alcanzadas.\footnote{%
    En la versión básica del modelo neoclásico, hemos analizado la determinación del equilibrio del
mercado definiendo un salario de reserva y una decisión de participación en el mercado, así como
las elasticidades mashallianas y hicksianas, la primera como la suma del efecto sustitución
(representada por la segunda) y el efecto renta. La oferta de trabajo estática se estudiaba
analizando el rol de las decisiones intrafamiliares.
La teoría dinámica de la oferta de trabajo da un papel fundamental a la posibilidad de sustituir
consumo y ocio en el tiempo, lo que también permite distinguir entre cambios transitorios o
permanentes.
}

\subsubsection{Obtención de la oferta intertemporal de trabajo}
Desde una perspectiva dinámica:
\begin{lnum}
    \item Un consumidor toma sus decisiones a lo largo del horizonte que
representa su ciclo vital, una sucesión de periodos, $t$, que comienza en $O$ y termina en una fecha, denotada por $T$. Por lo que su función de utilidad será la suma de sus Funciones de Utilidad instantáneas:\footnote{%
    Cabe recordar que la Función de Utilidad intertemporal es aditivamente separable por lo que no permite tener en cuenta
la inercia del consumo pasado o persistencia de hábitos, que revelan los estudios empíricos.\\ Para considerar este efecto habría que incorporar la influencia del consumo pasado en la función
de utilidad del periodo actual.
}
    $$\sum_{t=0}^TU(C_t, L_t, t)$$
    \item Ausencia de decisiones de formación.\footnote{%
        La formación incrementa el capital humano de un individuo y eleva sus previsiones de ganancias salariales, por lo que el trade off debe considerar
ocio, tiempo de trabajo y tiempo dedicado a estudiar.
    }
    \item Los individuos ahorran.
    \item Para cada periodo, la dotación total de tiempo que el agente puede dedicar a ocio o a trabajo, se supone constante y normalizada a 1. Las horas trabajadas durante un periodo $t$ son igual a ($1-L_t$). 
    \item Los consumidores perciben también otros ingresos. Por lo que la restricción presupuestaria que enfrenta el agente puede ser representado como:
    $$A_t=(1+r_t)A_{t-1}+B_t+w_t(1-L_t)-C_t$$
\end{lnum}  

El problema del consumidor es maximizar su utilidad intertemporal sujeto a la restricción presupuestaria descrita. Para resolverlo, se plantea la función auxiliar de Lagrange del problema
de maximización del consumidor.

\eqblock{
\begin{aligned}
    &\max_{\{C_t, L_t\}}{\sum_{t=0}^TU(C_t,L_t,t)} \qquad \text{s.a.}\quad (1+r_t)A_{t-1}+B_t+w_t(1-L_t)-C_t\leq A_t\quad \forall t\\[1em]
    \text{Donde,} \quad
    &\begin{cases}
        A_t\equiv \;\text{Activos del agente}\\[0.5em]
        B_t\equiv \;\text{Renta no salarial}\\[0.5em]
        R_t=B_t+r_tA_{t-1}\;\equiv \;\text{Renta no salarial total en $t$}\\
    \end{cases}\\[1em]
    \Longrightarrow &\mathcal{L}=\sum_{t=0}^TU(C_t,L_t,t)-\sum_{t=0}^Tv_t[A_t-(1+r_t)A_{t-1}-B_t-w_t(1-L_t)+C_t]\quad \text{\scriptsize ; ($v_t\equiv$ Multiplicador de Lagrange)}\\[2em]
    &\text{CPO's:}\\
    &\begin{cases}
        \left.\begin{aligned}
            &\text{(1)}\quad U_C(C_t,L_t,T)=v_t\\[0.5em]
            &\text{(2)}\quad U_L(C_t,L_t,T)=v_t\;w_t
        \end{aligned}\right\}\quad
        \begin{cases}
            \text{\scriptsize Óptimo Intratemporal (1 $\circ$ 2): $\frac{U_c}{U_L}=w_t$ ($=$ Modelo estático)}\\
            \text{\scriptsize \textbf{Óptimo Intertemporal} (2 $\circ$ 3, fundamental):}
            \quad \begin{aligned}
                &\text{\scriptsize Para dos periodos de tiempo,}\\
                &\text{\scriptsize $\frac{\partial U}{\partial L_t}=v_t\;w_t$ y $\frac{\partial U}{\partial L_{t+1}}=v_t\;w_{t+1}$}\\
                &\text{\scriptsize ya que relaciona: $(v_t, v_{t+1})\sim r$}
            \end{aligned}
        \end{cases}\\[1em]
        \text{(3)}\quad v_t=(1+r_{t+1})\;v_{t+1}\to \text{\tiny Ecuación de Euler: $(UMg_L)_t=(UMg_L)_{t+1}$ (por el factor descuento). $\ln{v_t}=-\sum_{\tau=1}^{\tau=t}\ln{(1+r_\tau)+\ln{v_0}}$}\\[0.5em]
    \end{cases}\\[2em]
    \Longrightarrow &\text{Demandas óptimas de bienes y ocio:}\qquad 
    \underbrace{\boxed{C_t=C(w_t,v_t, t)\quad \text{y}\quad L_t=L(w_t, v_t, t)}}_{\text{Demandas de \authorfont{FRISCH}}}
\end{aligned}
}{Oferta neoclásica (enfoque dinámico). Problema del consumidor.}

Por tanto, como vemos, de la obtención de las CPO (y operando intertemporalmente entre ellas) podemos destacar:
\begin{lnum}
    \item Condición intratemporal entre ocio y consumo ($\frac{U_L}{U_C}=w_t$). Al igual que en el caso estático, el equilibrio deberá garantizar que $RMS_C^L=w_t$ .
    \item Condición intertemporal de consumo ($U_C(\cdot, t)=(1+r_{t+1})U_C(\cdot, t+1)$) (o condición de Euler en el Consumo). Nos garantiza que en el óptimo, la utilidad de consumir una unidad hoy, debe ser igual a la utilidad de ahorrar esa unidad y consumir la mañana; y,
    \item Condición intertemporal del ocio ($\frac{U_L(\cdot, t)}{U_L(\cdot, t+1)}=(1+r_{t+1})\frac{w_t}{w_{t+1}}$) nos indica que la demanda relativa de ocio, y por tanto la oferta relativa de trabajo, responde a los salarios relativos. Será la relevante a la hora de valorar los efectos de shocks permanentes y transitorios en las decisiones de trabajo
\end{lnum}

Gracias a esta tercera conclusción podemos analizar cómo actuaría el individuo ante diferentes shocks. Si suponemos que el $ES>ER$\footnote{%
    Lo que extraemos de la conclusión de las CPO's es el Efecto Sustitución, por lo que para que debemos asumir que el primero domina al segundo. No se pueden dar conclusiones sin conocer cuál es la función de utilidad
}:
\begin{la}
    \item \textit{Shock transitorio}. Ante un aumento del salario en el momento actual, que deja inalterado el salario mañana el agente aprovechará el shock para trabajar más hoy en relación a mañana, incrementando por tanto su oferta de trabajo ($\uparrow\frac{w_t}{w_t+1}\;\to \;\frac{U_L(\cdot, t)}{U_L(\cdot, t+1)}\uparrow$).
    \item \textit{Shock permanente}. Ante un aumento del salario en $t$ y $t+1$, no habrá cambio en la oferta relativa de trabajo.
\end{la}

Este enfoque intertemporal de la oferta de trabajo es entonces crucial, porque nos permite ver cómo se modifican las decisiones de oferta de trabajo relativa ante cambios en los salarios relativos o en el tipo de interés. Esto es lo que LUCAS y RAPPING denominaron sustitución intertemporal de la oferta de trabajo.\footnote{%
    Este fenómeno resulta crucial en la Teoría del Ciclo Real iniciada por KYDLAND y PRESCOTT, por lo que esta relevancia ha justificado que se le haya prestado una especial atención en la literatura. Más concretamente, este fenómeno serviría para explicar por qué a lo largo del ciclo hay variaciones tan importantes en el empleo, mientras que los salarios reales apenas experimentan variaciones. Esta idea fue utilizada por KYDLAND y PRESCOTT para explicar que la economía está sometida a shocks que tienen repercusiones en la remuneración del trabajo y del capital, y que los agentes actúan ajustando sus decisiones de manera óptima a estos shocks, lo que les llevaría a alterar sus ofertas relativas de trabajo  y a producir dicho fenómeno de sustitución intertemporal de la oferta de trabajo. Sin embargo, esto precisaba que la elasticidad de sustitución intertemporal fuera muy elevada, mientras que la evidencia empírica en trabajos como el de HALL o el de BLUNDELL demuestran que dicha elasticidad es considerablemente reducida.
} 

\begin{specialblock}{Estudio de caso: Aplicaciones de la Oferta de trabajo dinámica: jubilación y pensiones}
    Muchos países de la OCDE tienen sistemas de pensiones de jubilación, públicas o privadas, que permiten a los trabajadores recibir una renta una vez que se han retirado del mercado de trabajo. Estos sistemas crean incentivos a los trabajadores para jubilarse antes o después. El modelo de ciclo de vida\footnote{%
        En el trabajo clásico de STOCK y WISE (1990) se hace una revisión exhaustiva de las principales cuestiones en torno a esta modelo (\textit{option value model}) y alguno de los modelos principales. Además, en el trabajo de BELLINI (2008) también se presentan los principales modelos al respecto.
    } considera una persona empleada en el periodo $\tau$ que decide retirarse en el periodo $s\ge \tau$. Suponemos que el individuo no trabaja después de haberse retirado. Llamamos $V_\tau(s)$ al bienestar del consumidor cuando decide retirarse en la fecha $s$, y $T_m$, a la edad legal de jubilación. El individuo elige su momento de jubilación, $s^*$, resolviendo el siguiente problema de maximización: $$\max_{s}{V_\tau(s)}\quad \text{s.a.}\quad T_m\ge s \ge \tau$$ La forma en que estimemos el momento en que se produce la jubilación nos permite conocer el valor de la opción real de no tomar la decisión en ese momento. Al existir incertidumbre en este proceso, retrasar la jubilación y tener la posibilidad de llevarla a cabo después tiene un valor que se traduce en un coste de oportunidad de jubilarse. En la medida en que este coste sea mayor a cero conviene seguir trabajando, pero si es menor que cero conviene tomar la decisión de jubilarse. Formalmente, si la decisión es irreversible, es decir, $s^*=\tau$, el individuo decide retirarse. Si $s^*>\tau$ el individuo continuaría trabajando y reconsideraría su decisión en el momento $(\tau+1)$. El valor de la opción de no retirarse en el momento actual sería igual a $V_\tau(s^*)-V_\tau(\tau)$. El valor de la opción de jubilarse se asociaría a la diferencia entre la utilidad esperada de jubilarse de forma inmediata y la utilidad esperada de hacerlo cuando exista la mayor utilidad esperada.\footnote{%
        GUSTMARI et al. (1994) observó que son aquellos individuos con pensiones más altas los que tienen incentivos para retirarse más pronto del mercado de trabajo, mientras que aquellos que tenían más presión financiera solían posponer esta decisión y alargar su vida laboral. Es posible por razones de eficiencia que las propias empresas ofrezcan planes de jubilación anticipada, y por ello, estas empresas suelen atraer trabajadores con una preferencia por la jubilación anticipada. Para eliminar este posible sesgo en la selección de trabajadores, existen muchos trabajos (LUMSDAINE et al, 1990) que también analizan el comportamiento de los trabajadores ante cambios anticipados en las decisiones de jubilación.
    }
\end{specialblock}

\subsubsection{Obtención de la demanda intertemporal de trabajo}
Al analizar la demanda de trabajo desde el punto de vista dinámico se suele adoptar una perspectiva distinta a la que se torna en el análisis dinámico de la oferta de trabajo. 

La empresa no sustituye intertemporalrnente trabajadores, a menos que haya costes de ajuste de su fuerza de trabajo.\textbf{Incluir artículo The Economist}\footnote{%
    Si no existen estos costes de ajuste en la cantidad de trabajo la empresa ajusta óptimamente en cada periodo su plantilla de manera que maximice beneficios.
} Por tanto, en este sentido el análisis de la demanda de trabajo es similar al de la demanda de capital. Las fluctuaciones de la demanda de trabajo se explicarían como respuestas a las perturbaciones que pudieran afectar a los ingresos de la empresa o a la productividad. En la medida en que el ajuste de plantillas conlleva costes, es necesario introducirlos en el análisis para poder dinamizar la demanda de trabajo.\footnote{%
    La teoría dinámica de la demanda de trabajo pennite analizar estas cuestiones y con ello la forma y la velocidad de los ajustes del empleo.
} El problema de la empresa es determinar la cantidad óptima de factores, y entre ellos el empleo, que maximice el valor presente de sus beneficios:
\eqblock{
\begin{aligned}
    &\max_{N(t)}{\int_0^\infty e^{-\phi t}\{p(t)F(N(t), t)-w(t)N(t)\}\mathrm{d}t}\quad \text{Donde,} 
    \begin{cases}
        F(N(t)) \equiv \text{Función de producción, con fp $L_t$}\\
        \phi \equiv \text{Factor descuento temporal}\\
        \text{Para simplificar: } N \equiv \text{Empleo total}
    \end{cases}\\[1em]
    &\text{De la CPO, con respecto a } N:\\[1em]
    &\hspace{6em}p(t\frac{\partial F}{\partial N}=w(t))\\[2em]
    &\text{\scriptsize Si suponemos que $F(N(t), t)=y(t)=A_t(K_t,N_t)$ ($A\equiv$ la tecnología),}\\[2em]
    &\text{\scriptsize en términos logarítmicos (elasticidades o cambios porcentuales), tenemos: $\ln{A_t}=a_t=a_0+\theta_t\quad \theta\sim(0,Var<\infty)$. Entonces,}\\
    \Longrightarrow &\text{La demanda de trabajo:} \qquad \hspace{6em}\boxed{n_t^d=n(k_t, a_t, \frac{w_t}{p_t})}\\[3em]
    &\text{\scriptsize Positiva en $k_t, a_t$ y negativa en $w_t/p_t$}
\end{aligned}
}{Enfoque dinámico. Demanda de trabajo y elasticidades (cambio porcentual)}

Si existe una perturbación positiva, la demanda de trabajo aumentaría, generando un aumento del empleo y del salario real, pero sólo temporal. El proceso se repetiría en cada periodo.\footnote{%
    El problema de este enfoque es que es necesario tener en cuenta la existencia de costes de ajuste en las decisiones sobre la mano de obra para tener una perspectiva completa de la dinámica del empleo. Sería por tanto más completo si estudiamos el impacto que tienen los costes de ajuste en la detenninación de la demanda de trabajo. Para ello, conviene acudir a otro enfoque, que vamos a analizar en temas posteriores, siguiendo el trabajo de NICKELL (1986).
}

Por tanto, como vemos, enriquecer el análisis neoclásico básico a través de este problema de optimización intertemporal nos permite definir un concepto de suma importancia para la macroeconomía moderna como es la sustitución intertemporal de la oferta de trabajo. Pero lo cierto es que el enfoque neoclásico básico presentado al inicio de la exposición puede enriquecerse añadiendo otro aspecto adicional, y es que los trabajadores en la realidad no son homogéneos. Pueden diferenciarse, por ejemplo, en su formación. 

\clearpage
\section{Teoría del capital humano}
\puntosclave{
Las diferencias salariales se pueden explicar por diferencias en la productividad de los
trabajadores, en la medida en que la adquisición de competencias y conocimientos permite
obtener un mejor salario.

Si el beneficio total obtenido en el mercado de trabajo posterior a educarse supera los costes
directos e indirectos invertidos durante el periodo de formación sería razonable pensar que
conviene invertir en educación.

La educación es un medio para acumular capital humano que ofrecería rendimiento futuro.
Dicha educación se ve influida también por las características del individuo y por su
habilidad innata.
}
En la práctica los trabajadores no son homogéneos y como tal no reciben una remuneración única. 

Los trabajadores pueden diferenciarse entre sí en diferentes aspectos, pero fundamentalmente cobra importancia las diferencias educativas o de formación de capital humano. A grandes rasgos, la literatura ha proporcionado dos explicaciones a por qué la educación puede llevar a los trabajadores a percibir salarios diferentes:\footnote{%
    De hecho, no son alternativos, sino complementarios, ya que uno es relativo a la formación de las decisiones de los individuos y la otra se centra en el análisis general de la educación como inversión. Frente al enfoque del capital humano, en el que las decisiones de educación constituyen una
inversión; en presencia de información asimétrica, la educación serviría para seleccionar
trabajadores. Desde esta perspectiva, el sistema de
educación juega un papel esencial de filtro: selecciona a los individuos basándose en su habilidad
intrínseca, permitiéndoles señalizar sus habilidades a los potenciales empleadores. Esto se
fundamenta en el hecho de que aquellos individuos que actúan de forma más efectiva en la vida
son aquellos que también actúan de forma más efectiva en los estudios, por lo que el nivel de
educación es una señal que permite a la empresa inferir el nivel de productividad del
trabajador que va a contratar
}
\begin{lnum}
    \item \textit{Educación como filtro}. Desarrollada en el trabajo seminal de SPENCE, defiende que la educación es un ejercicio de señalización en un marco de información asimétrica.\footnote{%
        [\authorfont{Ver Tema 3.A.13}] Las empresas, antes de contratar a un trabajador, desconocen cuál es la productividad de éste y, por tanto, son incapaces a priori de pagar salarios distintos a los agentes muy productivos de los poco productivos. La educación puede servir como señal, de manera que los agentes productivos se educan para transmitir creíblemente esta información a la empresa. La clave está en que la educación es más costosa para los no productivos que para los productivos, por lo que si un agente está dispuesto a educarse mucho es porque es productivo, por lo que la empresa usa dicha señal para poder diferenciar entre uno y otro. Esta teoría sostiene, por tanto, que la educación no mejora la productividad, sino que simplemente permite separar a los agentes según dicha productividad.
    }
    \item \textit{Teoría del capital humano.} Este enfoque fue iniciado por BECKER (1964).
\end{lnum}

Vamos a centrarnos en la Teoría del Capital Humano ya que las decisiones de inversión en educación son fundamentales para entender los diferenciales salariales del mercado de trabajo. El concepto de capital humano surge al intentar explicar que las personas gastan en sí mismas en determinadas ocasiones, no para tener satisfacciones presentes, sino para tener un beneficio futuro o potencial, tanto pecuniario como no pecuniario. Uno de ellos es el beneficio en términos de empleo y salario que esperan obtener si realizan una inversión que supone adquirir conocimiento y que aumenta la productividad, lo que acaba traduciéndose en una mayor renta. Por ello, en la medida en que los individuos tengan un nivel mayor de cualificación, su remuneración en el mercado de trabajo también se incrementa. La hipótesis de la teoría del capital humano supone que:\footnote{%
    La inversión en educación puede explicarse como la inversión en cualquier otro bien, que tiene en cuenta el flujo actualizado de beneficios. Éste es el enfoque de MCCONNELL y BRUE (1995).
} 

\textit{En la medida en que las diferencias salariales se pueden explicar por diferencias en la productividad de los trabajadores, la adquisición de competencias y conocimientos permite obtener un mejor salario.}

El coste de invertir en educación viene representado por los gastos directos que supone el período de estudios (matriculas, costes de viajes, compra de materiales, etc.), así como las pérdidas potenciales de renta debido al coste de oportunidad de no estar dedicado a una actividad remunerada y los costes psicológicos que surgen del estrés y la dificultad de estudiar. Si la inversión en educación da lugar a una acumulación de competencias, que eleva la productividad, proporcionará rendimientos en forma de una mayor remuneración futura. Podemos representar esto gráficamente:

\imagenfit{Coste-Beneficio_Khumano.png}{Curva Coste-Beneficio de la inversión en capital humano. \authorfont{McCONELL y BRUE}, 1997}

Básicamente, habría que comparar los beneficios brutos derivados de la inversión de educación, representados en el gráfico como el diferencial de ganancias de las dos curvas, con las áreas que representan los costes directos o indirectos. En la medida en que el diferencial de beneficio total bruto obtenido en el mercado de trabajo posterior al invertir en educación supere los costes directos e indirectos invertidos durante el periodo de fonnación, sería razonable pensar que conviene invertir en educación. 

\subsection{Aportaciones de Becker} 

\subsubsection{Tipos de educación}
Sobre esta base, cabe diferencias entre dos tipos de educación:\footnote{%
    Ya que no siempre la inversión en capital humano y, consecuentemente, una mejora en la productividad del asalariado conduce a un incremento de su salario.\\ Un trabajador que ha adquirido competencias y experiencia que mejoran su productividad sólo conseguirá un mayor pago si se enfrenta a dos o más empleadores que compitan por contratarle. De hecho, un único empleador no tendría incentivos para elevar el salario de un trabajador cuya productividad hubiera aumentado gracias a la inversión en educación, si ese trabajador no puede conseguir un trabajo mejor en otra parte. Esta observación llevó a BECKER a adoptar la distinción entre formación general, que potencia la productividad de un individuo para cualquier tipo de trabajo, y la formación específica, que sólo potencia la productividad en un único tipo de trabajo.
} 
\begin{la}
    \item \textit{La formación general}. asociada fundamentalmente al trabajador, que puede hacerla rentable en diferentes tipos de trabajo; y,
    \item \textit{La formación específica}. Asociada con un tipo particular de trabajo y, por tanto, con la empresa. Por ejemplo, es frecuente que las empresas de consultoría realicen un curso inicial a los trabajadores para explicar los principales procedimientos en el trabajo que éstas utilizan. 
\end{la}

Nos centraremos el estudio de las decisiones de inversión en capital humano relativo al primer tipo de formación, puesto que se trata de habilidades que pueden ser utilizadas en cualquier sector y cualquier empresa y, por tanto, existiendo competencia en el mercado, si la empresa no remunera adecuadamente la productividad del trabajador, éste podrá marcharse a otra empresa que sí le pague según el capital humano adquirido. 

\subsubsection{Modelo de \authorfont{BECKER}}
Siendo la educación un medio para acumular capital humano que generaría rendimiento futuro, lo primero que debemos señalar es que los años de educación se verán enfluenciados también por características del individuo y por su capital humano heredado.\footnote{%
    Podemos plantear que el individuo pueda elegir entre adquirir educación o trabajar y obtener una renta laboral.
} Una vez definamos las características implícitas del individuo, podremos formular su decisión de invertir en capital humano como:

\eqblock{
\begin{aligned}
    &h(t)=\theta\;\sigma(t)\;h(t)\qquad \text{Donde,}\quad
    \begin{cases}
        \text{\scriptsize$h(t)\equiv$ \; Designa el stock de capital humano, en $t$}\\[0.5em]
        \text{\scriptsize$\theta\equiv \;$ Parámetro de eficiencia del individuo}\\[0.5em]
        \text{\scriptsize$\sigma(t)\equiv \; \substack{\text{Esfuerzo en la educación, fracción de $t$}\\ \text{que éste dedica a formarse}}$}
    \end{cases}\\[1em]
    &\text{Entonces, el beneficio a lo largo del ciclo vital ($\Omega$):}\\
    &\hspace{5em}\Omega=\int_0^\tau [1-\sigma(t)]Ah(t)\;e^{-r\;t}\mathrm{d}t\\
    &\text{\scriptsize Donde, $Ah(t)\equiv\;x_i$ producidos en la fecha $t$ por $i$, que tiene un capital humano $h(t)$ }\\[2em]
    &\Longrightarrow \text{Por tanto,}\\
    &\frac{\partial \Omega}{\partial \sigma(t)}=-Ah(t)\;e^{-r\tau}+\int_t^\tau\theta\;A[1-\sigma(z)]h(z)\;e^{-rt}\mathrm{d}z \quad \text{y, siempre que $\underbrace{r>\theta}_{\substack{\text{\scriptsize Eficiencia $\theta$ sea}\\[0.2em]\text{\scriptsize $\ge r,$ rendimiento de $K$}\\}}$, entonces}\\[1em]
    &\hspace{5em}\frac{\partial}{\partial t}\left[\frac{\partial \Omega}{\partial \sigma(t)}\right]=Ah(t)\;e^{-r\;t}(r-\theta)\quad \Longrightarrow\quad \frac{(PMg_{\sigma(t)})}{t}>0
\end{aligned}
}{Modelo de \authorfont{BECKER}. Desarrollo analítico}

Esto significa que los gastos vinculados con la educación tienen que proporcionar al menos los mismos rendimientos que las inversiones financieras que podría realizar el agente que decide si invertir en educación:
\begin{la}
    \item Si $r<\theta$. La tasa a la que se acumula capital humano es superior a la tasa de descuento, entonces el rendimiento marginal de una inversión en educación hecha en cualquier fecha, $t>O$. es inferior al rendimiento marginal de la misma inversión hecha en la fecha inicial, $t=O$. Esto implicaría que el rendimiento marginal del esfuerzo en educación a lo largo de la vida del individuo sería decreciente. Bajo estas circunstancias, un agente siempre tiene interés en concentrar su educación en un periodo completo al comienzo de su vida. Por tanto, una persona recibiría educación si es suficientemente paciente para que los rendimientos de esa educación sean suficientemente altos; por el contrario,
    \item Si $\theta <r$. La tasa a la que se acumula capital humano es inferior a la tasa de descuento, entonces el rendimiento marginal de la inversión en educación en positivo y el agente enfrentará un problema de determinación de los años óptimos de formación.
\end{la}

\paragraph{Duración óptima de la formación}
La duración óptima de tiempo que un agente debería dedicar a adquirir educación ($S(\theta)$) al comienzo de su vida durante un intervalo de tiempo podríamos calcularla como:

\eqblock{
\begin{aligned}
    &S(\theta)=
    \begin{cases}
        T+\frac{1}{r}\ln{(\frac{\theta-r}{\theta})}, &\qquad \text{si} \quad \theta\ge \frac{r}{1-e^{-(r\;T)}}\\[0.5em]
        0&\qquad \text{en cualquier otro caso.}
    \end{cases}\\[2em]
    \Longrightarrow\text{Por tanto,} \quad 
    &\begin{cases}
        \uparrow T \quad& \\[0.5em]
        \uparrow \theta \quad&\to\qquad  \uparrow S(\theta)\\[0.5em]
        \uparrow \mathrm{gap}(\theta, r,T)=\theta(1-e^{-(rT)})-r&
    \end{cases}\hspace{1em}
    \begin{aligned}
        &\text{\scriptsize Esperanza de vida}\\[0.5em]
        &\text{\scriptsize Eficiencia innata}\\[0.5em]
        &\substack{\text{\scriptsize Eficiencia y jubilación/esperanza}\\\text{\scriptsize suficientemente grandes}}
    \end{aligned}
\end{aligned}
}{Modelo de \authorfont{BECKER}. Duración óptima de adquirir educación}

\paragraph{Valoración y extensiones}
La teoría del capital humano permite, por tanto, adentrarse en numerosos aspectos sobre las decisiones individuales con respecto a la educación. Existen autores que han demostrado que la teoría de capital humano explica la relación entre el salario y las ganancias derivadas del trabajo.\footnote{%
    El salario $w(\theta)$ aumenta con la eficiencia $\theta$ de la inversión en educación por dos razones. Por un lado, cada periodo de educación aumenta el stock de capital humano a un mayor nivel cuanto más eficiente es el individuo. Por otro lado, los individuos más eficientes suelen estudiar durante más tiempo. También se puede observar que el salario depende del stock inicial de conocimientos $h_0\sim h(t_0)$. En este sentido, el capital humano heredado también influye en la renta derivada del trabajo.
} Posteriormente, han existido numerosas extensiones a esta teoría:
\begin{lnum}
    \item HECKMAN (1976) elaboró un modelo tratando de analizar el trade-off entre el consumo y el ocio para la evidencia comentada (vid. supra);
    \item El trabajo de SPENCE (1973) sobre el capital humano como un modo de señalizar la productividad de un individuo dado la limitación de que la habilidad del individuo no sea una variable observable;
    \item \authorfont{BEN-PORATH} ( 1967) asume que el rendimiento marginal al esfuerzo en educación es decreciente a lo largo de la vida del individuo dado que hay una tasa de depreciación del conocimiento, $\delta$;\footnote{%
        Esta función asume, por tanto, que la acumulación de capital humano es una función cóncava del esfuerzo. El propósito de esta hipótesis es obtener las soluciones para las que, en cada periodo, el individuo pasaría parte de su tiempo recibiendo formación y parte trabajando.
    } y, 
    \item \authorfont{CAHUC et al.} (2014) representan el esfuerzo dedicado a la educación (gráfico de la izquierda), el stock de capital humano a lo largo de la vida (línea discontinua del gráfico de la derecha) y las ganancias salariales (línea continua del gráfico de la derecha)\footnote{%
        Estos gráficos representan la duración de educación y la duración de las ganancias salariales para los titulados universitarios en Estados Unidos. Se observa que la educación se acumula principalmente al inicio del ciclo vital, y que conforme se avanza a partir de una determinada edad, el esfuerzo en educación es menor. En cambio, el rendimiento salarial esperado es mayor en la medida en que se reduce el tiempo dedicado a estudiar y se sustituye por tiempo dedicado a trabajar. En la medida en que también el capital humano es mayor, el salario también remunera por una parte este esfuerzo. Si se modificará el parámetro de eficiencia (los autores partes de $\theta=0.4$), que refleja la heterogeneidad de las habilidades de los individuos, se podría representar también el perfil de los individuos con estudios secundarios en el mismo país, observando las diferencias
    }
\end{lnum}

\imagenfit{CAHUC-EEUU.png}{Evidencia empírica, EE.UU. \authorfont{CAHUC et. al.}, 2014}{Apuntes ICEX-CECO. Tema 3.A.25}

Ahora bien, la teoría del capital humano también está sujeta a críticas que hay que considerar: 
\begin{lnum}
    \item Si según lo visto el mercado es capaz de proveer el nivel óptimo de formación, ¿por qué se financia públicamente el capital humano? La teoría de BECKER no tiene en cuenta ni las externalidades ni las restricciones al crédito.\footnote{%
        Lo primero justificaría la necesidad de potenciar estas actividades educativas, ya que generan una externalidad positiva que no es tenida en cuenta por los individuos a la hora de decidir cuánto educarse. Las segundas indican que individuos que querrían educarse durante $S(\theta^*)$ años no pueden hacerlo porque no logran acceder a créditos que financien los costes de tal educación, por lo que el sector público debe tratar de ayudar a estos agentes a alcanzar el nivel educativo óptimo.
    }
    \item El modelo presupone que mayor educación lleva a una mayor productividad y, por tanto, es de esperar que lleve a mayores salarios. Pero ¿realmente existe esta dirección de causalidad? Esto ha llevado a una extensa cantidad de trabajos empíricos que han tratado de determinar la relación causal subyacente entre $\uparrow S(\theta)\to PMg_L\uparrow\;\to w\uparrow$ o $\uparrow PMg_L\to S(\theta)\uparrow\;\to w\uparrow$, pero resulta complejo obtener causalidad en este aspecto.\footnote{%
        Démonos cuenta que Spence argumenta la dirección contraria: me educo más porque soy más productivo y eso me lleva a tener mayores salarios.\\ Quizás uno de los más interesantes es el desarrollado por DUFLO, quien se apoya en un masivo programa de construcción de escuelas en Indonesia para estimar el impacto causal de la educación en los salarios, encontrando que cada año adicional de escolarización permitía a los individuos acceder a salarios un 10\% superiores. Según este trabajo de DUFLO, por tanto, la relación causal sería la propuesta por GARY BECKER: mayor educación mejora la productividad y permite obtener mayores salarios.
    } 
\end{lnum}

\subsection{Ecuaciones de ingresos de Mincer}
Una de las ecuaciones más conocidas en la estimación de los rendimientos educativos es la de MINCER (1974). Viene definida como: 

\eqblock{
\begin{aligned}
    &\log{w}=\beta_0+\beta_1+\beta_2\;\mathrm{Exp}+\beta_3\;\mathrm{Exp}^2+\epsilon\\[1em]
    \text{Donde,}
    &\begin{cases}
        w&\equiv \text{Ingresos del individuo (incremental, log)}\\[0.5em]
        s&\equiv \text{Número de años de educación}\\[0.5em]
        \mathrm{Exp}&\equiv \text{Número de años de experiencia}
    \end{cases}\\[1em]
\end{aligned}
}{Modelo del Capital Humano. Ecuación de ingresos de \authorfont{MINCER}}

La variable de interés sería $\beta_1$, que permite calcular cuál sería el cambio porcentual en las ganancias resultante de un año adicional de educación\footnote{%
    Lo que se define como la tasa de retorno de la educación o el rendimiento de dicha educación. \\ Por su parte, los coeficientes de la experiencia se interpretan como medidas del impacto de la formación en el empleo sobre los salarios.
}. La tasa de crecimiento de las ganancias disminuye con el tiempo porque los trabajadores acumulan menos capital humano cuando se hacen mayores.

\imagenfit{MINCER.png}{Curva de regresión sobre ganancias ($w$), en función de la edad}{Apuntes ICEX-CECO. Tema 3.A.25}

La popularidad de esta ecuación se ha extendido a diversos contextos y países dada la facilidad para poder generar un resultado comprensible y cuantificable del efecto de la educación y la experiencia laboral. No obstante, con el tiempo también se vio que existían numerosas limitaciones al estimar esta ecuación por mínimos cuadrados ordinarios (MCO).\footnote{%
    Algunas de estas limitaciones son la existencia de variables omitidas, la suposición de que la educación se trata como una variable exógena (cuando en algunos casos es endógena), la existencia de sesgos de selección o la existencia de varias tasas de retorno de la educación (no una sola).
}


\clearpage
\section{Economía empírica}
\puntosclave{
A la hora de estimar con datos la oferta de trabajo existen algunas limitaciones metodológicas principalmente centradas en los errores de medida y en los problemas de selección.
 
Por medio de experimentos naturales se puede identificar de forma causal y rigurosa el efecto de distintas políticas públicas en el mercado de trabajo.

Dos ejemplos de este tipo de evidencia son los referentes a los salarios mínimos y su impacto en el empleo y al capital humano y los rendimientos de la educación.
}

\subsection{Regularidades empíricas sobre la oferta de trabajo}


\subsection{Experimentos naturales}










\clearpage
\section{Conclusión}

\subsection{Recapitulación}

En esta exposición hemos presentado los rasgos más destacables de la teoría neoclásica del mercado de trabajo, análisis al que hemos añadido un enfoque intertemporal y la posibilidad de invertir en capital humano para enriquecer el enfoque. 

En este sentido, es indudable que el marco neoclásico es un marco excesivamente simplificado que lleva a la irreal conclusión de que no existe desempleo involuntario. El mercado de trabajo presenta agentes con poder de mercado, tal y como analizan los modelos de negociación como el modelo de Derecho a Gestionar de NICKEL y ANDREWS. También es un mercado que presenta heterogeneidad entre trabajadores e información asimétrica que hace que trabajadores y empresas formen parte de un proceso de matching que requiere tiempo para encasillarlos, dando lugar al paro friccional. También puede estar sujeto a problemas de riesgo moral, como pone de manifiesto el modelo de Saphiro y STIGLITZ.

No obstante, esto no supone de ninguna forma que el marco neoclásico carezca de valor. Debemos tomarlo como un baseline, en el que bajo unos supuestos idóneos pero irreales nos permiten alcanzar unos resultados sumamente interesantes, y es la inexistencia del desempleo involuntario. Levantar los supuestos restrictivos tal y como hacen los trabajos mencionados nos permite ver cómo alejarnos de los mismos puede modificar enormemente las conclusiones alcanzadas.

Finalmente, me gustaría terminar recalcando la importancia que cobra la evidencia empírica en los temas relativos al mercado de trabajo. Debemos darnos cuenta que las conclusiones alcanzadas ante una determinada cuestión pueden modificarse de manera notoria según la perspectiva desde la que lo analicemos.

\subsection{Extensiones y relación con otras partes del Temario}



\subsection{Opinión}

\subsubsection{Idea final (Salida o cierra)}

\clearpage
\pagenumbering{gobble}
\MyBibliography
\MyBibItem{Autor}{Título}{Año}{DOI -- LINK}


% --- Anexos (no aparecen en el índice) ---
\clearpage
\anexo{\textbf{Preguntas Test}}
%\input{\TemaCode/questions.tex} 


\end{document}